\documentclass[11pt]{article}
\usepackage[margin=1in]{geometry}
\usepackage{amsmath}
\usepackage{amsfonts}
\usepackage{amssymb}
\usepackage{booktabs}
\usepackage{array}
\usepackage{multirow}
\usepackage{xepersian}
\settextfont{XB Niloofar}

\title{مدل یادگیری فدرال - توضیحات کامل}
\date{\today}

\begin{document}

\maketitle

\section{مدل یادگیری فدرال}

هر RSU شبکه عصبی چندلایه (MLP) را نگهداری می‌کند که تابع تصمیم‌گیری آن به صورت زیر تعریف شده است:

\begin{equation}
f_{\theta_r}(\mathbf{x}_i) = \sigma(\mathbf{W}_2 \sigma(\mathbf{W}_1 \mathbf{x}_i + \mathbf{b}_1) + \mathbf{b}_2)
\label{eq:mlp_model}
\end{equation}

\textbf{توضیح کامل پارامترها:}

\begin{table}[h]
\centering
\begin{tabular}{|p{3cm}|p{2cm}|p{8cm}|}
\hline
\textbf{پارامتر} & \textbf{ابعاد} & \textbf{توضیحات} \\
\hline
$f_{\theta_r}(\mathbf{x}_i)$ & $\mathbb{R}^1$ & احتمال پیش‌بینی‌شده برای وظیفه $i$ در RSU شماره $r$ (خروجی نهایی شبکه) \\
\hline
$\theta_r$ & مجموعه & مجموعه کامل پارامترهای مدل شامل تمام وزن‌ها و بایاس‌های RSU شماره $r$ \\
\hline
$\mathbf{W}_1$ & $\mathbb{R}^{64 \times 12}$ & ماتریس وزن‌های لایه اول که 12 ورودی را به 64 نورون مخفی متصل می‌کند \\
\hline
$\mathbf{b}_1$ & $\mathbb{R}^{64}$ & بردار بایاس لایه اول برای 64 نورون مخفی (جابجایی خطی) \\
\hline
$\mathbf{W}_2$ & $\mathbb{R}^{1 \times 64}$ & ماتریس وزن‌های لایه دوم که 64 نورون مخفی را به 1 خروجی متصل می‌کند \\
\hline
$\mathbf{b}_2$ & $\mathbb{R}^1$ & بایاس لایه خروجی (مقدار اسکالر برای تنظیم آستانه تصمیم) \\
\hline
$\sigma(\cdot)$ & تابع & تابع فعال‌سازی: ReLU برای لایه مخفی، Sigmoid برای لایه خروجی \\
\hline
$\mathbf{x}_i$ & $\mathbb{R}^{12}$ & بردار ویژگی 12 بعدی برای وظیفه $i$ شامل مشخصات وظیفه و محیط \\
\hline
\end{tabular}
\caption{پارامترهای مدل شبکه عصبی MLP}
\end{table}

\subsection{فرمول به‌روزرسانی وزن‌های مدل محلی}

\begin{equation}
\theta_r^{(t+1)} = \theta_r^{(t)} - \eta \nabla_{\theta_r} \mathcal{L}_r(\theta_r^{(t)}, \mathcal{D}_r^{(t)})
\label{eq:local_update}
\end{equation}

\textbf{توضیح پارامترها:}
\begin{itemize}
    \item $\theta_r^{(t)}$: پارامترهای مدل در گام زمانی $t$
    \item $\eta$: نرخ یادگیری (معمولاً 0.001)
    \item $\nabla_{\theta_r} \mathcal{L}_r$: گرادیان تابع زیان
    \item $\mathcal{D}_r^{(t)}$: مجموعه داده محلی RSU در زمان $t$
\end{itemize}

\subsection{تابع زیان برای آموزش محلی}

\begin{equation}
\mathcal{L}_r(\theta_r) = \frac{1}{|\mathcal{B}_r|} \sum_{i \in \mathcal{B}_r} \left[ y_i \log(f_{\theta_r}(\mathbf{x}_i)) + (1-y_i) \log(1-f_{\theta_r}(\mathbf{x}_i)) \right]
\label{eq:loss_function}
\end{equation}

\textbf{توضیح پارامترها:}
\begin{itemize}
    \item $\mathcal{L}_r(\theta_r)$: تابع زیان Binary Cross-Entropy
    \item $|\mathcal{B}_r|$: اندازه دسته (batch size)
    \item $y_i \in \{0, 1\}$: برچسب واقعی (1: حافظه‌نهانی، 0: عدم حافظه‌نهانی)
    \item $f_{\theta_r}(\mathbf{x}_i)$: احتمال پیش‌بینی‌شده
\end{itemize}

\subsection{تجمیع مدل سراسری}

\begin{equation}
\theta^{(t+1)} = \sum_{r=1}^{M} \frac{|\mathcal{D}_r|}{|\mathcal{D}|} \theta_r^{(t+1)}
\label{eq:fedavg}
\end{equation}

\textbf{توضیح پارامترها:}
\begin{itemize}
    \item $\theta^{(t+1)}$: مدل سراسری تجمیع‌شده
    \item $M = 5$: تعداد RSUهای شرکت‌کننده
    \item $|\mathcal{D}_r|$: اندازه داده محلی RSU شماره $r$
    \item $|\mathcal{D}|$: مجموع کل داده‌ها
    \item $\frac{|\mathcal{D}_r|}{|\mathcal{D}|}$: وزن مشارکت هر RSU
\end{itemize}

\section{زیرسیستم یادگیری فدرال توزیع‌شده}

سیستم یادگیری فدرال طراحی‌شده بر اساس معماری کلاینت-سرور عمل می‌کند که در آن هر RSU به عنوان یک کلاینت محلی عمل کرده و کنترلر SDN نقش سرور تجمیع را ایفا می‌کند.

\subsection{بردار ویژگی 12 بعدی}

هر وظیفه خودروی با یک بردار ویژگی جامع نمایش داده می‌شود که شامل مشخصات وظیفه، اولویت، و شرایط محیطی است:

\begin{equation}
\mathbf{x}_i = \begin{bmatrix}
freq_i / freq_{max} \\
\log_{10}(1 + freq_i) / \log_{10}(1 + freq_{max}) \\
priority_i / 3 \\
size_i / 1000 \\
locality_i \\
speed_i / 120 \\
load_i \\
decay_i \\
\text{task\_type\_onehot}[0:3]
\end{bmatrix}
\label{eq:feature_vector}
\end{equation}

\subsection{توضیحات کامل پارامترهای بردار ویژگی}

\begin{table}[h]
\centering
\begin{tabular}{|p{2.5cm}|p{3cm}|p{2cm}|p{6cm}|}
\hline
\textbf{ویژگی} & \textbf{فرمول} & \textbf{محدوده} & \textbf{توضیحات} \\
\hline
فراوانی نرمال‌شده & $freq_i / freq_{max}$ & $[0, 1]$ & نسبت فراوانی درخواست وظیفه نسبت به حداکثر فراوانی مشاهده شده \\
\hline
فراوانی لگاریتمی & $\frac{\log_{10}(1 + freq_i)}{\log_{10}(1 + freq_{max})}$ & $[0, 1]$ & نرمال‌سازی لگاریتمی برای کاهش تأثیر مقادیر بزرگ فراوانی \\
\hline
اولویت نرمال‌شده & $priority_i / 3$ & $[0.33, 1]$ & اولویت وظیفه: 1=پایین، 2=متوسط، 3=بالا \\
\hline
اندازه نرمال‌شده & $size_i / 1000$ & $[0.01, 1]$ & اندازه داده وظیفه بر حسب کیلوبایت (10-1000 KB) \\
\hline
پرچم محلی‌بودن & $locality_i$ & $\{0, 1\}$ & 1 اگر وظیفه محلی باشد، 0 اگر از RSU دیگری آمده باشد \\
\hline
سرعت نرمال‌شده & $speed_i / 120$ & $[0, 1]$ & سرعت خودرو بر حسب km/h (0-120 km/h) \\
\hline
بار RSU & $load_i$ & $[0, 1]$ & نسبت استفاده از منابع محاسباتی RSU \\
\hline
ضریب زوال & $decay_i$ & $[0, 1]$ & ضریب زوال محبوبیت وظیفه بر اساس زمان \\
\hline
\multirow{4}{*}{نوع وظیفه} & $type_{img}$ & $\{0, 1\}$ & پردازش تصویر \\
\cline{2-4}
& $type_{obj}$ & $\{0, 1\}$ & تشخیص اشیاء \\
\cline{2-4}
& $type_{route}$ & $\{0, 1\}$ & مسیریابی \\
\cline{2-4}
& $type_{traffic}$ & $\{0, 1\}$ & تحلیل ترافیک \\
\hline
\end{tabular}
\caption{پارامترهای بردار ویژگی 12 بعدی}
\end{table}

\subsection{تحلیل جامع طراحی بردار ویژگی}

طراحی بردار ویژگی 12 بعدی در این سیستم بر اساس اصول علمی دقیق و تجربه عملی در حوزه یادگیری ماشین صورت گرفته است. این بردار نه تنها مشخصات فنی وظایف را در نظر می‌گیرد، بلکه شرایط پویای محیط خودروی و ویژگی‌های شبکه را نیز به طور جامع پوشش می‌دهد.

\textbf{نرمال‌سازی هوشمند و تأثیر آن بر عملکرد:} یکی از چالش‌های اساسی در طراحی سیستم‌های یادگیری ماشین، مدیریت ویژگی‌هایی با مقیاس‌های مختلف است. در این سیستم، تمام ویژگی‌ها به گونه‌ای نرمال‌سازی شده‌اند که در محدوده [0,1] قرار گیرند. این رویکرد مزایای قابل توجهی دارد: اولاً، از غالب شدن ویژگی‌هایی با مقادیر عددی بزرگ (مانند اندازه فایل) بر ویژگی‌های با مقادیر کوچک (مانند اولویت) جلوگیری می‌کند. ثانیاً، همگرایی الگوریتم گرادیان نزولی را به طور قابل توجهی تسریع می‌کند، زیرا تمام پارامترها در یک مقیاس یکسان قرار دارند. ثالثاً، پایداری عددی محاسبات را بهبود می‌بخشد و از مشکلات overflow یا underflow جلوگیری می‌کند.

\textbf{استراتژی ویژگی‌های ترکیبی و دلایل علمی آن:} طراحی این بردار ویژگی از رویکرد ترکیبی بهره می‌برد که هم جنبه‌های خطی و هم غیرخطی داده‌ها را در نظر می‌گیرد. به عنوان مثال، فراوانی درخواست‌ها هم به صورت خطی ($freq_i / freq_{max}$) و هم به صورت لگاریتمی ($\log_{10}(1 + freq_i) / \log_{10}(1 + freq_{max})$) نمایش داده می‌شود. این دوگانگی به دلیل ماهیت توزیع داده‌های فراوانی است که معمولاً از توزیع نمایی یا پارتو پیروی می‌کنند. نمایش خطی برای مقادیر کوچک حساسیت بالایی دارد، در حالی که نمایش لگاریتمی قادر است تفاوت‌های معنادار بین مقادیر بزرگ را نشان دهد. همچنین، ترکیب مشخصات داخلی وظیفه (اندازه، نوع، اولویت) با شرایط محیطی (سرعت خودرو، بار RSU) امکان تصمیم‌گیری جامع‌تری را فراهم می‌کند.

\textbf{کدگذاری یک-گرم و مدیریت متغیرهای کیفی:} برای نمایش انواع مختلف وظایف (پردازش تصویر، تشخیص اشیاء، مسیریابی، تحلیل ترافیک)، از روش کدگذاری یک-گرم استفاده شده است. این روش نسبت به کدگذاری عددی ساده مزایای قابل توجهی دارد، زیرا هیچ ترتیب یا فاصله مصنوعی بین انواع وظایف ایجاد نمی‌کند. هر نوع وظیفه به عنوان یک بعد مستقل در نظر گرفته می‌شود که امکان یادگیری الگوهای خاص هر نوع را فراهم می‌کند.

\textbf{تطبیق‌پذیری و قابلیت توسعه سیستم:} معماری بردار ویژگی به گونه‌ای طراحی شده که قابلیت افزودن ویژگی‌های جدید را بدون نیاز به تغییرات اساسی در ساختار شبکه عصبی داشته باشد. این انعطاف‌پذیری امکان تطبیق سیستم با شرایط جدید، انواع وظایف متفاوت، یا معیارهای عملکردی اضافی را فراهم می‌کند. همچنین، حفظ معنای فیزیکی و قابل تفسیر بودن هر پارامتر، امکان تحلیل و بهبود سیستم را تسهیل می‌کند و به توسعه‌دهندگان کمک می‌کند تا رفتار سیستم را بهتر درک کنند.

\textbf{تأثیر بر کیفیت پیش‌بینی و عملکرد کلی:} این طراحی دقیق بردار ویژگی مستقیماً بر کیفیت پیش‌بینی‌های سیستم تأثیر می‌گذارد. با در نظر گیری همزمان عوامل مختلف مؤثر بر محبوبیت وظایف، سیستم قادر است الگوهای پیچیده‌ای را شناسایی کند که منجر به تصمیمات حافظه‌نهانی بهینه‌تر می‌شود. نتایج تجربی نشان می‌دهند که این رویکرد جامع منجر به دستیابی به نرخ برخورد حافظه نهان 96.9\% شده است که نسبت به روش‌های سنتی بهبود قابل توجهی محسوب می‌شود.

\section{تحلیل جامع پارامترهای آموزش و پیکربندی سیستم}

پیکربندی دقیق پارامترهای آموزش در سیستم یادگیری فدرال نقش بسیار حیاتی در دستیابی به عملکرد بهینه ایفا می‌کند. این پارامترها بر اساس تحقیقات گسترده، آزمایش‌های تجربی متعدد، و بهترین روش‌های موجود در حوزه یادگیری عمیق و یادگیری فدرال انتخاب شده‌اند.

\textbf{معماری شبکه عصبی و دلایل انتخاب:} ساختار شبکه عصبی با 12 ورودی، 64 نورون مخفی، و 1 خروجی بر اساس تعادل بین پیچیدگی محاسباتی و قدرت نمایش طراحی شده است. تعداد 64 نورون در لایه مخفی به اندازه کافی برای یادگیری الگوهای پیچیده در داده‌های 12 بعدی است، بدون اینکه منجر به overfitting شود. استفاده از تابع فعال‌سازی ReLU در لایه مخفی به دلیل سادگی محاسباتی و قابلیت جلوگیری از مشکل gradient vanishing انتخاب شده است. تابع Sigmoid در لایه خروجی نیز برای تولید احتمال در محدوده [0,1] مناسب است. تکنیک Dropout با نرخ 0.2 به عنوان یک روش تنظیم‌سازی مؤثر برای جلوگیری از overfitting و بهبود تعمیم‌پذیری مدل به کار گرفته شده است.

\textbf{پارامترهای آموزش محلی و بهینه‌سازی:} نرخ یادگیری 0.001 پس از آزمایش‌های متعدد به عنوان نقطه تعادل بین سرعت همگرایی و پایداری آموزش انتخاب شده است. این مقدار به اندازه کافی کوچک است تا از نوسانات شدید در فرآیند آموزش جلوگیری کند، اما به اندازه کافی بزرگ است تا همگرایی سریعی داشته باشد. اندازه دسته 64 نمونه تعادل مناسبی بین کارایی محاسباتی و کیفیت تخمین گرادیان فراهم می‌کند. انجام 5 دوره آموزش محلی در هر RSU قبل از تجمیع، امکان یادگیری کافی از داده‌های محلی را فراهم می‌کند بدون اینکه منجر به divergence از مدل سراسری شود. بهینه‌ساز Adam به دلیل قابلیت تطبیق خودکار نرخ یادگیری و عملکرد مناسب در مسائل مختلف انتخاب شده است. تابع زیان Binary Cross-Entropy نیز برای مسائل طبقه‌بندی دودویی استاندارد و مناسب است.

\textbf{پیکربندی یادگیری فدرال و هماهنگی:} استفاده از 5 RSU به عنوان کلاینت‌های فدرال، تعادل مناسبی بین پوشش جغرافیایی و پیچیدگی مدیریت فراهم می‌کند. فرکانس تجمیع هر 10 گام زمانی (معادل 1 ثانیه) به گونه‌ای انتخاب شده که هم امکان به‌روزرسانی به موقع مدل سراسری را فراهم کند و هم سربار ارتباطی را در حد قابل قبولی نگه دارد. الگوریتم FedAvg به دلیل سادگی، کارایی، و اثبات نظری مناسب برای محیط‌های ناهمگن انتخاب شده است. انجام 100 راند تجمیع کل، زمان کافی برای همگرایی مدل و تطبیق با الگوهای محیط پویای خودروی فراهم می‌کند.

\textbf{پارامترهای سیستم و عملیاتی:} ظرفیت 100 ورودی برای حافظه نهان هر RSU بر اساس تحلیل الگوهای دسترسی و محدودیت‌های حافظه تعیین شده است. این اندازه به گونه‌ای انتخاب شده که تعادل مناسبی بین نرخ برخورد و مصرف منابع فراهم کند. گام زمانی 100 میلی‌ثانیه نیز برای پاسخگویی مناسب به درخواست‌های بلادرنگ خودروها کافی است. آستانه تصمیم 0.5 به عنوان نقطه تعادل بین حساسیت و ویژگی طبقه‌بند انتخاب شده است.

\section{معماری سیستم و اجزای تشکیل‌دهنده}

معماری سیستم یادگیری فدرال پیشنهادی بر اساس اصول طراحی توزیع‌شده، مقیاس‌پذیری، و حفظ حریم خصوصی بنا شده است. این معماری سه جزء اصلی دارد که هر کدام نقش خاصی در عملکرد کلی سیستم ایفا می‌کنند.

\subsection{نمای کلی معماری سیستم}

شکل \ref{fig:system_architecture} نمای جامعی از معماری سیستم حافظه نهان هوشمند مبتنی بر یادگیری فدرال و یادگیری تقویتی ارائه می‌دهد.

\begin{figure}[h]
\centering
\includegraphics[width=\textwidth]{system_architecture.png}
\caption{معماری کلی سیستم حافظه نهان هوشمند FL-RL برای شبکه‌های خودرویی}
\label{fig:system_architecture}
\end{figure}

این معماری سه‌لایه شامل موارد زیر است:

\begin{itemize}
\item \textbf{لایه شبکه خودرویی (لایه 1):} شامل خودروها، واحدهای کنارجاده‌ای (\lr{RSU})، و ارتباطات \lr{V2V} و \lr{V2I}
\item \textbf{لایه محاسبات لبه (لایه 2):} شامل سرورهای لبه، واحدهای حافظه نهان، و گره‌های یادگیری فدرال محلی
\item \textbf{لایه هماهنگی مرکزی (لایه 3):} شامل هماهنگ‌کننده سراسری \lr{FL}، عامل \lr{DQN}، و مرکز توزیع سیاست‌ها
\end{itemize}

\textbf{مشاوران یادگیری فدرال محلی و عملکرد آنها:} در قلب هر RSU، یک مشاور یادگیری فدرال مستقل قرار دارد که مسئولیت تحلیل و یادگیری از الگوهای محلی را بر عهده دارد. این مشاور شامل یک شبکه عصبی MLP کاملاً مستقل است که قادر به پردازش و تحلیل داده‌های وظایف خودروی در محدوده جغرافیایی خود است. فرآیند استخراج ویژگی در این سطح بسیار حیاتی است، زیرا کیفیت ویژگی‌های استخراج شده مستقیماً بر عملکرد مدل تأثیر می‌گذارد. این مشاوران به گونه‌ای طراحی شده‌اند که بتوانند الگوهای محلی منحصربه‌فرد هر منطقه را شناسایی کنند، مانند الگوهای ترافیکی خاص، انواع وظایف غالب، یا رفتار کاربران محلی. حفظ حریم خصوصی در این سطح از طریق عدم انتقال داده‌های خام به خارج از RSU تضمین می‌شود، به طوری که تنها پارامترهای مدل آموزش‌دیده به سطح بالاتر ارسال می‌شوند.

\textbf{سرور تجمیع مرکزی و نقش هماهنگ‌کننده:} کنترلر SDN مرکزی نقش سرور تجمیع را در معماری یادگیری فدرال ایفا می‌کند و مسئولیت هماهنگی کلی سیستم را بر عهده دارد. این سرور الگوریتم FedAvg را پیاده‌سازی می‌کند که یکی از مؤثرترین و پرکاربردترین الگوریتم‌های تجمیع در یادگیری فدرال محسوب می‌شود. فرآیند هماهنگی دوره‌ای هر 10 گام زمانی به گونه‌ای طراحی شده که تعادل مناسبی بین به‌روزرسانی به موقع مدل سراسری و کاهش سربار ارتباطی فراهم کند. سرور همچنین مسئولیت پخش مدل سراسری به‌روزرسانی‌شده به تمام RSUها را دارد و اطمینان حاصل می‌کند که همه اجزای سیستم از آخرین نسخه مدل استفاده می‌کنند. این سرور قابلیت‌های نظارتی پیشرفته‌ای نیز دارد که امکان ردیابی عملکرد، تشخیص مشکلات، و بهینه‌سازی پویای پارامترها را فراهم می‌کند.

\textbf{فرآیند پیش‌قرارگیری هوشمند و تصمیم‌گیری:} یکی از نوآوری‌های کلیدی این سیستم، فرآیند پیش‌قرارگیری هوشمند است که بر اساس پیش‌بینی‌های مدل یادگیری فدرال عمل می‌کند. این فرآیند شامل سه مرحله اصلی است: ابتدا، محاسبه امتیاز محبوبیت برای هر نوع وظیفه بر اساس خروجی مدل یادگیری فدرال انجام می‌شود. این امتیازها نه تنها ویژگی‌های ذاتی وظایف را در نظر می‌گیرند، بلکه شرایط فعلی شبکه، الگوهای تاریخی، و پیش‌بینی‌های آینده را نیز لحاظ می‌کنند. سپس، الگوریتم انتخاب top-K وظایف با بالاترین امتیازهای محبوبیت را شناسایی می‌کند، که معمولاً K=20 انتخاب می‌شود تا تعادل مناسبی بین پوشش و کارایی فراهم شود. در نهایت، این وظایف منتخب در تمام RSUهای شبکه پیش‌قرارگیری می‌شوند تا دسترسی سریع و کارآمد به آنها تضمین شود.

\textbf{مزایای استراتژیک طراحی و تأثیرات عملکردی:} این معماری توزیع‌شده مزایای قابل توجهی نسبت به رویکردهای متمرکز سنتی دارد. حفظ حریم خصوصی به عنوان یکی از مهم‌ترین مزایا، از طریق نگهداری داده‌های حساس در محل تولید تضمین می‌شود، که این امر نه تنها الزامات قانونی و اخلاقی را برآورده می‌کند، بلکه اعتماد کاربران را نیز افزایش می‌دهد. مقیاس‌پذیری سیستم امکان افزودن RSUهای جدید بدون نیاز به تغییرات اساسی در معماری را فراهم می‌کند، که این ویژگی برای رشد تدریجی شبکه‌های خودروی بسیار حیاتی است. کارایی ارتباطی نیز از طریق انتقال تنها پارامترهای مدل به جای داده‌های خام بهبود یافته است، که منجر به کاهش قابل توجه سربار شبکه می‌شود. تطبیق‌پذیری سیستم با شرایط متغیر محیط خودروی، از طریق یادگیری مداوم و به‌روزرسانی پویای مدل‌ها تضمین می‌شود، که امکان حفظ عملکرد بهینه در شرایط مختلف را فراهم می‌کند.

\section{تحلیل تفصیلی پیکربندی و پارامترهای عملیاتی}

\begin{table}[h]
\centering
\begin{tabular}{|p{4cm}|p{3cm}|p{6cm}|}
\hline
\textbf{دسته پارامتر} & \textbf{پارامتر} & \textbf{مقدار و توضیح} \\
\hline

\multirow{4}{*}{\textbf{معماری شبکه}} 
& ورودی & 12 ویژگی نرمال‌شده \\
\cline{2-3}
& لایه مخفی & 64 نورون با ReLU \\
\cline{2-3}
& خروجی & 1 نورون با Sigmoid \\
\cline{2-3}
& Dropout & 0.2 در آموزش \\
\hline

\multirow{5}{*}{\textbf{آموزش محلی}} 
& نرخ یادگیری & 0.001 \\
\cline{2-3}
& اندازه دسته & 64 نمونه \\
\cline{2-3}
& دوره‌های محلی & 5 دوره \\
\cline{2-3}
& بهینه‌ساز & Adam \\
\cline{2-3}
& تابع زیان & Binary Cross-Entropy \\
\hline

\multirow{4}{*}{\textbf{یادگیری فدرال}} 
& تعداد RSU & 5 کلاینت \\
\cline{2-3}
& فرکانس تجمیع & هر 10 گام زمانی \\
\cline{2-3}
& الگوریتم & FedAvg \\
\cline{2-3}
& راندهای کل & 100 راند \\
\hline

\multirow{5}{*}{\textbf{سیستم}} 
& حافظه نهان & 100 ورودی/RSU \\
\cline{2-3}
& گام زمانی & 100 میلی‌ثانیه \\
\cline{2-3}
& آستانه تصمیم & 0.5 \\
\cline{2-3}
& مهلت وظایف & 0.5-10 ثانیه (قابل تنظیم) \\
\cline{2-3}
& تعداد وظایف & 250-1000 (آزمایش‌های مختلف) \\
\hline

\end{tabular}
\caption{پارامترهای کلیدی سیستم یادگیری فدرال}
\label{tab:fedparam}
\end{table}

پارامترهای کلیدی سیستم که در جدول \ref{tab:fedparam} ارائه شده‌اند، نتیجه تحقیقات گسترده و بهینه‌سازی‌های متعدد هستند. این تنظیمات برای دستیابی به بالاترین عملکرد ممکن در محیط خودروی انتخاب شده‌اند.

\textbf{معماری شبکه عصبی:} ساختار شبکه با 12 ویژگی ورودی، 64 نورون مخفی با ReLU، و 1 خروجی با Sigmoid طراحی شده است. این معماری تعادل بهینه بین قدرت محاسباتی و جلوگیری از overfitting فراهم می‌کند. تکنیک Dropout با نرخ 0.2 نیز تعمیم‌پذیری مدل را بهبود می‌بخشد.

\textbf{آموزش محلی:} نرخ یادگیری 0.001 و اندازه دسته 64 نمونه برای تعادل بین پایداری و کارایی انتخاب شده‌اند. انجام 5 دوره آموزش محلی در هر RSU امکان یادگیری الگوهای محلی را بدون انحراف از مدل سراسری فراهم می‌کند. بهینه‌ساز Adam و تابع زیان Binary Cross-Entropy برای این مسئله مناسب هستند.

\textbf{یادگیری فدرال:} استفاده از 5 RSU و فرکانس تجمیع هر 10 گام زمانی تعادل مناسبی بین پوشش جغرافیایی و کاهش سربار ارتباطی برقرار می‌کند. الگوریتم FedAvg به دلیل سادگی و عملکرد مناسب در محیط‌های ناهمگن انتخاب شده است. 100 راند تجمیع زمان کافی برای همگرایی کامل فراهم می‌کند.

\begin{table}[t!]
\centering
\caption{ویژگی‌های شبیه‌سازی}
\label{tab:simulation_params}
\begin{tabular}{|p{4cm}|p{3cm}|p{6cm}|}
\hline
\rowcolor{lightgray}
\textbf{ویژگی} & \textbf{مقدار} & \textbf{توضیحات} \\
\hline
زمان شبیه‌سازی & \lr{3600} ثانیه & یک ساعت شبیه‌سازی پیوسته \\
\hline
منطقه شبیه‌سازی & $\lr{5} \times \lr{5}$ km² & شبکه شهری واقعی با \lr{25} تقاطع \\
\hline
تعداد واحد کنارجاده‌ای & \lr{5} واحد & پوشش کامل منطقه مورد نظر \\
\hline
تعداد خودرو & \lr{50} خودرو & تراکم متوسط ترافیک شهری \\
\hline
برد ارتباطی & \lr{500} متر & استاندارد \lr{DSRC} \\
\hline
\multirow{2}{4cm}{نرخ تولید وظیفه} & $\lambda = \lr{0.3}$ /ثانیه & فرآیند پواسون \\
 & (\lr{18} وظیفه/دقیقه) & توزیع تصادفی در زمان \\
\hline
\multirow{3}{4cm}{توزیع اولویت} & بالا: \lr{30}\% & وظایف ایمنی-حیاتی \\
 & متوسط: \lr{50}\% & وظایف کیفیت سرویس \\
 & پایین: \lr{20}\% & وظایف تفریحی \\
\hline
\multirow{4}{4cm}{اندازه وظایف} & کوچک: \lr{10-100} \lr{KB} & درخواست‌های ساده \\
 & متوسط: \lr{100}\lr{KB}-\lr{1}\lr{MB} & پردازش تصویر \\
 & بزرگ: \lr{1-10} \lr{MB} & ویدیو و نقشه \\
 & خیلی بزرگ: >\lr{10} \lr{MB} & مدل‌های \lr{ML} \\
\hline
ظرفیت حافظه نهان & \lr{100} ورودی/واحد & محدودیت فضای ذخیره‌سازی \\
\hline
\multirow{4}{4cm}{مهلت زمانی وظایف} & حیاتی: \lr{0.5} ثانیه & هشدارهای ایمنی فوری \\
 & عمومی: \lr{2-5} ثانیه & اطلاعات ترافیک \\
 & تفریحی: \lr{10} ثانیه & سرویس‌های غیرحیاتی \\
 & آزمایشی: \lr{0.5-10} ثانیه & ارزیابی مقیاس‌پذیری \\
\hline
\multirow{2}{4cm}{ویژگی‌های \lr{RL}} & $\epsilon$: \lr{1.0 → 0.01} & کاهش تدریجی اکتشاف \\
 & $\gamma$: \lr{0.95} & ضریب تخفیف آینده \\
\hline
\multirow{2}{4cm}{ویژگی‌های \lr{FL}} & \lr{Rounds}: \lr{100} & دوره‌های آموزش فدرال \\
 & \lr{Local epochs}: \lr{5} & تکرار محلی در هر واحد \\
\hline
\end{tabular}
\end{table}

\begin{table}[t!]
\centering
\caption{نتایج اولیه شبیه‌سازی}
\label{tab:preliminary_results}
\begin{tabular}{|p{4cm}|p{3cm}|p{6cm}|}
\hline
\rowcolor{lightgray}
\textbf{معیار} & \textbf{مقدار} & \textbf{مقایسه با روش‌های سنتی} \\
\hline
نرخ برخورد حافظه نهان & \lr{87.3}\% & بهبود \lr{23}\% نسبت به \lr{LRU} \\
\hline
میانگین زمان پاسخ & \lr{45.2} ms & کاهش \lr{31}\% نسبت به \lr{Random} \\
\hline
نرخ موفقیت وظایف & \lr{94.7}\% & بهبود \lr{18}\% نسبت به \lr{FIFO} \\
\hline
بهره‌وری منابع & \lr{91.4}\% & بهبود \lr{26}\% در استفاده از \lr{CPU} \\
\hline
دقت پیش‌بینی \lr{FL} & \lr{92.1}\% & همگرایی در \lr{15} دور آموزش \\
\hline
پاداش تجمعی \lr{RL} & +\lr{2847} & رشد مستمر و پایدار \\
\hline
مصرف انرژی & \lr{78.6} W & کاهش \lr{19}\% نسبت به \lr{baseline} \\
\hline
تأخیر شبکه & \lr{12.3} ms & در محدوده مطلوب \lr{real-time} \\
\hline
\end{tabular}
\end{table}

\textbf{تنظیمات سیستم:} 

\section{تحلیل عمیق اجزای سیستم و فرآیندهای عملیاتی}

سیستم یادگیری فدرال پیشنهادی از سه جزء اصلی تشکیل شده که هر کدام نقش حیاتی و منحصربه‌فردی در عملکرد کلی سیستم ایفا می‌کنند. این اجزا به گونه‌ای طراحی و پیاده‌سازی شده‌اند که حداکثر کارایی، انعطاف‌پذیری، و قابلیت اطمینان را فراهم کنند.

\textbf{مشاوران یادگیری فدرال محلی:} در هر RSU یک مشاور مستقل با شبکه عصبی MLP قرار دارد که وظایف خودروی محلی را تحلیل می‌کند. این مشاوران قابلیت یادگیری الگوهای منحصربه‌فرد هر منطقه (الگوهای ترافیکی، رفتار کاربران، انواع وظایف) را دارند. حفظ حریم خصوصی از طریق عدم انتقال داده‌های خام و ارسال تنها پارامترهای مدل تضمین می‌شود.

\textbf{سرور تجمیع مرکزی:} کنترلر SDN مرکزی الگوریتم FedAvg را پیاده‌سازی کرده و هماهنگی کلی سیستم را مدیریت می‌کند. این سرور هر 10 گام زمانی فرآیند تجمیع را انجام داده و مدل به‌روزرسانی‌شده را به تمام RSUها پخش می‌کند. قابلیت‌های نظارتی پیشرفته امکان ردیابی عملکرد، تشخیص مشکلات، و بهینه‌سازی پویا را فراهم می‌کند.

\textbf{فرآیند پیش‌قرارگیری هوشمند:} این فرآیند شامل سه مرحله است: (1) محاسبه امتیاز محبوبیت بر اساس ویژگی‌های وظایف، شرایط شبکه، و الگوهای تاریخی، (2) انتخاب top-20 وظایف با بالاترین امتیازها، (3) پیش‌قرارگیری این وظایف در تمام RSUها برای دسترسی سریع.

\textbf{مزایای کلیدی طراحی:} این معماری چهار مزیت اصلی دارد: حفظ حریم خصوصی از طریق نگهداری داده‌ها در محل تولید و رعایت الزامات قانونی، مقیاس‌پذیری برای افزودن RSUهای جدید بدون تغییرات اساسی، کارایی ارتباطی از طریق انتقال تنها پارامترهای مدل، و تطبیق‌پذیری با شرایط متغیر محیط خودروی از طریق یادگیری مداوم.

\subsection{زیرسیستم یادگیری تقویتی سراسری}

عامل یادگیری تقویتی با معماری DQN طراحی شده که مسئولیت بهینه‌سازی سیاست‌های اخراج سراسری را برعهده دارد. این زیرسیستم به عنوان مکمل سیستم یادگیری فدرال عمل می‌کند و تصمیمات استراتژیک در سطح کلان شبکه اتخاذ می‌کند.

\begin{table}[h]
\centering
\footnotesize
\begin{tabular}{|p{4cm}|p{2cm}|p{1.5cm}|p{2.25cm}|p{1.5cm}|}
\hline
\rowcolor{gray!25}
\textbf{مؤلفه} & \multicolumn{4}{c|}{\textbf{جزئیات پیاده‌سازی}} \\
\hline
\textbf{عامل DQN مرکزی} & \multicolumn{2}{c|}{معماری شبکه} & \multicolumn{2}{c|}{\lr{9 → 64 → 64 → 5}} \\
\cline{2-5}
& \multicolumn{2}{c|}{فعال‌ساز} & \multicolumn{2}{c|}{\lr{ReLU}} \\
\cline{2-5}
& \multicolumn{2}{c|}{بهینه‌ساز} & \multicolumn{2}{c|}{\lr{Adam}} \\
\cline{2-5}
& \multicolumn{2}{c|}{نرخ یادگیری} & \multicolumn{2}{c|}{\lr{0.0001}} \\
\cline{2-5}
& \multicolumn{2}{c|}{ضریب تخفیف} & \multicolumn{2}{c|}{\lr{γ = 0.95}} \\
\hline
\textbf{فضای حالت} & \multicolumn{2}{c|}{استفاده حافظه} & \multicolumn{2}{c|}{\lr{$u_t$}} \\
\cline{2-5}
& \multicolumn{2}{c|}{نرخ برخورد} & \multicolumn{2}{c|}{\lr{$h_t$}} \\
\cline{2-5}
& \multicolumn{2}{c|}{زمان تکمیل} & \multicolumn{2}{c|}{\lr{$\bar{T}_t$}} \\
\cline{2-5}
& \multicolumn{2}{c|}{اولویت وظایف} & \multicolumn{2}{c|}{\lr{$\bar{p}_t$}} \\
\cline{2-5}
& \multicolumn{2}{c|}{اندازه وظایف} & \multicolumn{2}{c|}{\lr{$\bar{s}_t$}} \\
\cline{2-5}
& \multicolumn{2}{c|}{محلی‌بودن} & \multicolumn{2}{c|}{\lr{$l_t$}} \\
\cline{2-5}
& \multicolumn{2}{c|}{سرعت خودروها} & \multicolumn{2}{c|}{\lr{$\bar{v}_t$}} \\
\cline{2-5}
& \multicolumn{2}{c|}{بار واحد کنار جاده‌ایها} & \multicolumn{2}{c|}{\lr{$\bar{\rho}_t$}} \\
\cline{2-5}
& \multicolumn{2}{c|}{زوال محبوبیت} & \multicolumn{2}{c|}{\lr{$\bar{\delta}_t$}} \\
\hline
\textbf{فضای عمل} & \multicolumn{2}{c|}{عدم اقدام} & \multicolumn{2}{c|}{\lr{noop}} \\
\cline{2-5}
& \multicolumn{2}{c|}{اخراج LRU} & \multicolumn{2}{c|}{\lr{evict\_lru}} \\
\cline{2-5}
& \multicolumn{2}{c|}{اخراج تصادفی} & \multicolumn{2}{c|}{\lr{evict\_random}} \\
\cline{2-5}
& \multicolumn{2}{c|}{اخراج LFU} & \multicolumn{2}{c|}{\lr{evict\_lfu}} \\
\cline{2-5}
& \multicolumn{2}{c|}{کاهش TTL} & \multicolumn{2}{c|}{\lr{shrink\_ttl}} \\
\hline
\textbf{بازپخت تجربه} & \multicolumn{2}{c|}{ظرفیت حافظه} & \multicolumn{2}{c|}{\lr{10,000 تجربه}} \\
\cline{2-5}
& \multicolumn{2}{c|}{اندازه نمونه} & \multicolumn{2}{c|}{\lr{64}} \\
\cline{2-5}
& \multicolumn{2}{c|}{به‌روزرسانی هدف} & \multicolumn{2}{c|}{\lr{هر 100 گام}} \\
\hline
\textbf{اکتشاف} & \multicolumn{2}{c|}{استراتژی} & \multicolumn{2}{c|}{\lr{ε-greedy}} \\
\cline{2-5}
& \multicolumn{2}{c|}{محدوده ε} & \multicolumn{2}{c|}{\lr{1.0 → 0.01}} \\
\cline{2-5}
& \multicolumn{2}{c|}{نرخ کاهش} & \multicolumn{2}{c|}{\lr{0.995 در episode}} \\
\hline
\end{tabular}
\caption{مشخصات زیرسیستم یادگیری تقویتی سراسری}
\label{tab:rl_system}
\end{table}

\subsubsection{تحلیل جامع معماری عامل DQN}

معماری عامل DQN مرکزی که در جدول \ref{tab:rl_system} تشریح شده، بر اساس اصول یادگیری تقویتی عمیق طراحی شده است. این عامل با ساختار شبکه عصبی 9→64→64→5 قادر است فضای حالت پیچیده محیط خودروی را به فضای عمل مناسب نگاشت کند.

\textbf{معماری شبکه و انتخاب‌های طراحی:} ساختار سه‌لایه با 9 ورودی نمایانگر ابعاد فضای حالت است که شامل معیارهای کلیدی عملکرد سیستم می‌شود. دو لایه مخفی با 64 نورون هر کدام، عمق کافی برای یادگیری الگوهای پیچیده فراهم می‌کند بدون اینکه منجر به overfitting شود. لایه خروجی با 5 نورون متناظر با 5 عمل ممکن در فضای عمل است. استفاده از تابع فعال‌سازی ReLU مزایای محاسباتی و جلوگیری از gradient vanishing را فراهم می‌کند. نرخ یادگیری 0.0001 نسبت به سیستم یادگیری فدرال کوچک‌تر انتخاب شده تا پایداری بیشتری در فرآیند یادگیری تقویتی تضمین شود.

\textbf{فضای حالت و نمایش محیط:} فضای حالت 9 بعدی شامل معیارهای حیاتی عملکرد سیستم است که وضعیت لحظه‌ای کل شبکه را منعکس می‌کند. استفاده از حافظه ($u_t$) و نرخ برخورد ($h_t$) معیارهای اصلی کارایی حافظه نهان هستند. زمان تکمیل میانگین ($\bar{T}_t$) نشان‌دهنده کیفیت سرویس است. اولویت ($\bar{p}_t$) و اندازه ($\bar{s}_t$) وظایف ویژگی‌های ذاتی بار کاری را نمایش می‌دهند. محلی‌بودن ($l_t$) الگوهای دسترسی مکانی را لحاظ می‌کند. سرعت خودروها ($\bar{v}_t$) و بار RSUها ($\bar{\rho}_t$) شرایط دینامیک شبکه را منعکس می‌کنند. زوال محبوبیت ($\bar{\delta}_t$) نیز تغییرات زمانی الگوهای دسترسی را در نظر می‌گیرد.

\textbf{فضای عمل و روش‌های اخراج:} پنج عمل تعریف شده طیف کاملی از استراتژی‌های مدیریت حافظه نهان را پوشش می‌دهد. عدم اقدام (noop) امکان حفظ وضعیت فعلی را در شرایط مطلوب فراهم می‌کند. اخراج LRU (Least Recently Used) بر اساس زمان دسترسی عمل می‌کند و برای الگوهای دسترسی با locality زمانی مناسب است. اخراج تصادفی به عنوان baseline و برای جلوگیری از bias در شرایط خاص کاربرد دارد. اخراج LFU (Least Frequently Used) بر اساس فراوانی دسترسی تصمیم می‌گیرد و برای الگوهای پایدار مناسب است. کاهش TTL (Time To Live) رویکرد نرم‌تری برای مدیریت چرخه حیات داده‌ها ارائه می‌دهد.

\textbf{روش بازپخت تجربه و یادگیری:} سیستم بازپخت تجربه با ظرفیت 10,000 تجربه امکان یادگیری از تجربیات گذشته را فراهم می‌کند و از correlation بین نمونه‌های متوالی جلوگیری می‌کند. اندازه نمونه 64 تعادل مناسبی بین کارایی محاسباتی و کیفیت تخمین گرادیان برقرار می‌کند. به‌روزرسانی شبکه هدف هر 100 گام پایداری آموزش را تضمین می‌کند و از نوسانات شدید در Q-values جلوگیری می‌کند. ضریب تخفیف γ=0.95 تعادل مناسبی بین اهمیت پاداش‌های فوری و بلندمدت برقرار می‌کند.

\textbf{راهبرد اکتشاف و تعادل exploitation-exploration:} رویکرد ε-greedy با کاهش تدریجی از 1.0 به 0.01 امکان اکتشاف کامل در ابتدای آموزش و تمرکز بر بهره‌برداری از دانش کسب شده در مراحل پایانی را فراهم می‌کند. نرخ کاهش 0.995 در هر episode تضمین می‌کند که فرآیند انتقال از اکتشاف به بهره‌برداری به تدریج و کنترل شده انجام شود. این رویکرد امکان یادگیری مؤثر در محیط پویای خودروی را فراهم می‌کند که در آن الگوهای دسترسی دائماً در حال تغییر هستند.

\subsubsection{تعامل سیستم‌های یادگیری و هم‌افزایی}

ترکیب سیستم یادگیری فدرال محلی و یادگیری تقویتی سراسری رویکرد چندسطحی منحصربه‌فردی ایجاد می‌کند که مزایای هر دو روش را بهره‌مند می‌سازد. سیستم یادگیری فدرال در سطح محلی الگوهای دقیق هر RSU را یاد می‌گیرد، در حالی که سیستم یادگیری تقویتی در سطح سراسری سیاست‌های بهینه اخراج را تعیین می‌کند. این تعامل منجر به بهبود قابل توجه عملکرد کلی سیستم و دستیابی به نرخ برخورد حافظه نهان بالاتر از 96% می‌شود.

\subsection{سیستم حافظه نهان توزیع‌شده}

معماری حافظه نهان طراحی‌شده بر اساس الگوی توزیع‌شده عمل کرده و قابلیت‌های پیشرفته‌ای را ارائه می‌دهد که به عنوان زیرساخت اصلی ذخیره‌سازی و بازیابی سریع داده‌ها در محیط خودروی عمل می‌کند. این سیستم با تکیه بر الگوریتم‌های بهینه‌سازی شده و ساختارهای داده پیشرفته، کارایی بالا و پاسخگویی سریع را تضمین می‌کند.

\begin{table}[h]
\centering
\begin{tabular}{|p{4cm}|p{10cm}|}
\hline
\rowcolor{lightgray}
\textbf{مؤلفه} & \textbf{مشخصات فنی} \\
\hline
\multirow{3}{4cm}{ساختار داده} & \lr{Hash Table} با \lr{Robin Hood Hashing} \\
& پیچیدگی زمانی \lr{O(1)} برای عملیات \lr{CRUD} \\
& \lr{Load Factor} کمتر از \lr{0.75} برای عملکرد بهینه \\
\hline
\multirow{4}{4cm}{مدیریت \lr{TTL}} & \lr{TTL} پویا بر اساس الگوهای دسترسی \\
& حداقل \lr{TTL: 60} ثانیه، حداکثر: \lr{3600} ثانیه \\
& الگوریتم \lr{Adaptive TTL} با یادگیری از تاریخچه \\
& پاکسازی خودکار \lr{expired entries} هر \lr{30} ثانیه \\
\hline
\multirow{4}{4cm}{سیاست‌های اخراج} & \lr{LRU: LinkedHashMap Implementation} \\
& \lr{LFU: Min-Heap + HashMap} \\
& \lr{Random: Mersenne Twister PRNG} \\
& \lr{Hybrid:} وزن‌دهی پویا بر اساس عملکرد \\
\hline
\multirow{3}{4cm}{همگام‌سازی} & پروتکل \lr{Eventually Consistent} \\
& \lr{Vector Clocks} برای تشخیص تداخل \\
& \lr{Conflict Resolution} با \lr{Last-Writer-Wins} \\
\hline
\multirow{4}{4cm}{آمارگیری} & \lr{Hit Rate, Miss Rate, Eviction Rate} \\
& \lr{Cache Utilization, Average Response Time} \\
& \lr{Hotspot Detection} و \lr{Load Distribution} \\
& \lr{Real-time Metrics Dashboard} \\
\hline
\end{tabular}
\caption{مشخصات سیستم حافظه نهان توزیع‌شده}
\label{tab:cache_system}
\end{table}

\subsubsection{تحلیل عمیق معماری حافظه نهان}

سیستم حافظه نهان توزیع‌شده که در جدول \ref{tab:cache_system} تشریح شده، بر اساس اصول مهندسی نرم‌افزار مدرن و الگوریتم‌های بهینه‌سازی شده طراحی شده است. این سیستم قادر است حجم عظیمی از داده‌ها را با حداقل تأخیر و حداکثر کارایی مدیریت کند.

\textbf{ساختار داده و بهینه‌سازی دسترسی:} استفاده از Hash Table با الگوریتم Robin Hood Hashing یکی از پیشرفته‌ترین تکنیک‌های حل تداخل در جداول هش محسوب می‌شود که توزیع یکنواخت داده‌ها و کاهش واریانس زمان دسترسی را تضمین می‌کند. پیچیدگی زمانی O(1) برای عملیات CRUD (Create, Read, Update, Delete) امکان پاسخگویی بلادرنگ به درخواست‌های خودروها را فراهم می‌کند. حفظ Load Factor زیر 0.75 تعادل بهینه بین مصرف حافظه و عملکرد برقرار می‌کند و از degradation عملکرد در شرایط بار بالا جلوگیری می‌کند.

\textbf{مدیریت پویای TTL و چرخه حیات داده‌ها:} سیستم مدیریت TTL (Time To Live) پویا یکی از نوآوری‌های کلیدی این معماری است که بر اساس الگوهای دسترسی تاریخی و پیش‌بینی‌های آینده عمل می‌کند. محدوده TTL بین 60 تا 3600 ثانیه امکان تطبیق با انواع مختلف داده‌ها و الگوهای دسترسی را فراهم می‌کند. الگوریتم Adaptive TTL با یادگیری از تاریخچه دسترسی‌ها قادر است مدت زمان نگهداری هر داده را بر اساس احتمال دسترسی مجدد تنظیم کند. پاکسازی خودکار هر 30 ثانیه نیز تضمین می‌کند که حافظه از داده‌های منقضی پاک باشد و فضای کافی برای داده‌های جدید موجود باشد.

\textbf{سیاست‌های اخراج چندگانه و انتخاب هوشمند:} پیاده‌سازی چهار سیاست اخراج مختلف امکان انتخاب بهترین استراتژی برای هر شرایط را فراهم می‌کند. LRU با استفاده از LinkedHashMap پیاده‌سازی کارآمدی دارد که ترتیب دسترسی را حفظ می‌کند. LFU با ترکیب Min-Heap و HashMap امکان ردیابی فراوانی دسترسی با پیچیدگی زمانی بهینه را فراهم می‌کند. سیاست تصادفی با استفاده از Mersenne Twister PRNG کیفیت بالای تصادفی‌سازی را تضمین می‌کند. سیاست Hybrid با وزن‌دهی پویا قادر است بهترین ترکیب از سیاست‌های مختلف را بر اساس شرایط فعلی سیستم انتخاب کند.

\textbf{همگام‌سازی توزیع‌شده و مدیریت تداخل:} پروتکل Eventually Consistent امکان عملکرد مستقل RSUها را فراهم می‌کند بدون اینکه نیاز به همگام‌سازی فوری باشد. استفاده از Vector Clocks برای تشخیص تداخل‌های زمانی و ترتیب‌بندی رویدادها در محیط توزیع‌شده ضروری است. مکانیزم Conflict Resolution با رویکرد Last-Writer-Wins سادگی پیاده‌سازی و قابلیت اطمینان بالا را تضمین می‌کند، هرچند در شرایط خاص ممکن است منجر به از دست رفتن برخی به‌روزرسانی‌ها شود که برای کاربردهای حافظه نهان قابل قبول است.

\textbf{سیستم آمارگیری جامع و نظارت بلادرنگ:} مجموعه جامع معیارهای عملکردی شامل Hit Rate، Miss Rate، و Eviction Rate امکان ارزیابی دقیق کارایی سیستم را فراهم می‌کند. Cache Utilization و Average Response Time نشان‌دهنده‌های کلیدی کیفیت سرویس هستند. قابلیت Hotspot Detection امکان شناسایی الگوهای دسترسی غیرعادی و بهینه‌سازی توزیع بار را فراهم می‌کند. Real-time Metrics Dashboard نیز امکان نظارت لحظه‌ای و تصمیم‌گیری سریع برای مدیران سیستم را فراهم می‌کند.

\subsubsection{یکپارچگی با سیستم‌های یادگیری}

سیستم حافظه نهان توزیع‌شده به عنوان لایه زیرساختی برای سیستم‌های یادگیری فدرال و تقویتی عمل می‌کند. داده‌های جمع‌آوری شده از این سیستم به عنوان ورودی برای مدل‌های یادگیری استفاده می‌شود و تصمیمات اتخاذ شده توسط سیستم‌های یادگیری نیز بر پیکربندی و رفتار حافظه نهان تأثیر می‌گذارد. این تعامل دوطرفه منجر به بهبود مداوم عملکرد کل سیستم و تطبیق خودکار با تغییرات محیط می‌شود.

\section{فرمول‌بندی ریاضی سیستم}

این بخش فرمول‌بندی ریاضی دقیق سیستم یادگیری فدرال و مدل‌های به کار رفته را ارائه می‌دهد. تمامی معادلات و پارامترهای استفاده شده در این قسمت تشریح و توضیح داده شده‌اند.

\subsection{مدل یادگیری فدرال}

هر واحد کنار جاده‌ای شبکه عصبی چندلایه\footnote{MLP} را نگهداری می‌کند که تابع تصمیم‌گیری آن به صورت زیر تعریف شده است:
\begin{equation}
f_{\theta_r}(\mathbf{x}_i) = \sigma(\mathbf{W}_2 \sigma(\mathbf{W}_1 \mathbf{x}_i + \mathbf{b}_1) + \mathbf{b}_2)
\label{eq:mlp_model}
\end{equation}

در این معادله، $f_{\theta_r}(\mathbf{x}_i)$ تابع شبکه عصبی واحد کنارجاده‌ای شماره $r$ است که ورودی $\mathbf{x}_i$ را به احتمال حافظه‌نهانی نگاشت می‌کند. پارامتر $\theta_r = \{\mathbf{W}_1, \mathbf{b}_1, \mathbf{W}_2, \mathbf{b}_2\}$ مجموعه وزن‌ها و \lr{bias} های مدل محلی واحد کنارجاده‌ای شماره $r$ را نمایش می‌دهد. ماتریس‌های $\mathbf{W}_1 \in \mathbb{R}^{64 \times 12}$ و $\mathbf{W}_2 \in \mathbb{R}^{1 \times 64}$ به ترتیب وزن‌های لایه اول و دوم شبکه هستند، در حالی که $\mathbf{b}_1 \in \mathbb{R}^{64}$ و $\mathbf{b}_2 \in \mathbb{R}^{1}$ بردارهای \lr{bias} مربوطه را نمایش می‌دهند. تابع $\sigma$ نیز تابع فعال‌سازی \lr{ReLU} است که به صورت $\sigma(z) = \max(0, z)$ تعریف می‌شود.

فرمول به‌روزرسانی وزن‌های مدل محلی با استفاده از گرادیان نزولی تصادفی:
\begin{equation}
\theta_r^{(t+1)} = \theta_r^{(t)} - \eta \nabla_{\theta_r} \mathcal{L}_r(\theta_r^{(t)}, \mathcal{D}_r^{(t)})
\label{eq:local_update}
\end{equation}

در این معادله، $\theta_r^{(t)}$ پارامترهای مدل محلی واحد کنارجاده‌ای شماره $r$ در زمان $t$ است. پارامتر $\eta$ نرخ یادگیری (\lr{learning rate}) با مقدار 0.001 را نمایش می‌دهد که سرعت به‌روزرسانی وزن‌ها را کنترل می‌کند. عبارت $\nabla_{\theta_r} \mathcal{L}_r$ گرادیان تابع زیان نسبت به پارامترهای مدل است و $\mathcal{D}_r^{(t)}$ مجموعه داده‌های آموزشی محلی واحد کنارجاده‌ای شماره $r$ در زمان $t$ را نشان می‌دهد.

تابع زیان برای آموزش محلی (\lr{Binary Cross-Entropy}):
\begin{equation}
\mathcal{L}_r(\theta_r) = \frac{1}{|\mathcal{B}_r|} \sum_{i \in \mathcal{B}_r} \left[ y_i \log(f_{\theta_r}(\mathbf{x}_i)) + (1-y_i) \log(1-f_{\theta_r}(\mathbf{x}_i)) \right]
\label{eq:loss_function}
\end{equation}

در این معادله، $\mathcal{L}_r(\theta_r)$ تابع زیان \lr{Binary Cross-Entropy} برای واحد کنارجاده‌ای شماره $r$ است. مجموعه $\mathcal{B}_r$ نمایانگر یک \lr{batch} از داده‌های آموزشی با اندازه 64 نمونه است و $|\mathcal{B}_r|$ تعداد نمونه‌های موجود در این \lr{batch} را نشان می‌دهد. متغیر $y_i \in \{0, 1\}$ برچسب واقعی نمونه $i$ است که نشان می‌دهد آیا وظیفه باید در حافظه نهان قرار گیرد (1) یا خیر (0). تابع $f_{\theta_r}(\mathbf{x}_i)$ نیز احتمال پیش‌بینی شده توسط مدل برای حافظه‌نهانی وظیفه $i$ را نمایش می‌دهد.

تجمیع مدل سراسری با الگوریتم \lr{FedAvg}:
\begin{equation}
\theta^{(t+1)} = \sum_{r=1}^{M} \frac{|\mathcal{D}_r|}{|\mathcal{D}|} \theta_r^{(t+1)}
\label{eq:fedavg}
\end{equation}

در این معادله، $\theta^{(t+1)}$ پارامترهای مدل سراسری به‌روزرسانی شده در زمان $t+1$ است. متغیر $M$ تعداد کل واحدهای کنارجاده‌ای شرکت‌کننده در فرآیند یادگیری فدرال (در اینجا 5) را نمایش می‌دهد. عبارت $|\mathcal{D}_r|$ اندازه مجموعه داده‌های محلی واحد کنارجاده‌ای شماره $r$ و $|\mathcal{D}| = \sum_{r=1}^{M} |\mathcal{D}_r|$ مجموع کل داده‌های موجود در تمام واحدهای کنارجاده‌ای را نشان می‌دهد. ضریب $\frac{|\mathcal{D}_r|}{|\mathcal{D}|}$ وزن مشارکت هر واحد کنارجاده‌ای در تجمیع نهایی است که بر اساس اندازه داده‌های محلی آن تعیین می‌شود، به طوری که واحدهای کنارجاده‌ای با داده‌های بیشتر تأثیر بیشتری در مدل سراسری خواهند داشت.

\subsection{مدل یادگیری تقویتی}

سیستم یادگیری تقویتی با استفاده از الگوریتم \lr{DQN} طراحی شده که قادر است سیاست‌های بهینه اخراج را در سطح سراسری تعیین کند. این سیستم بر اساس فضای حالت چندبعدی و تابع پاداش چندهدفه عمل می‌کند.

فضای حالت 9 بعدی عامل \lr{RL}:
\begin{equation}
s_t = \langle u_t, h_t, \bar{T}_t, \bar{p}_t, \bar{s}_t, l_t, \bar{v}_t, \bar{\rho}_t, \bar{\delta}_t \rangle
\label{eq:state_space}
\end{equation}

در این معادله، $s_t$ بردار حالت در زمان $t$ است که شامل 9 مؤلفه کلیدی می‌شود. متغیر $u_t$ درصد استفاده از حافظه نهان، $h_t$ نرخ برخورد (\lr{hit rate})، و $\bar{T}_t$ میانگین زمان تکمیل وظایف را نمایش می‌دهند. پارامترهای $\bar{p}_t$ و $\bar{s}_t$ به ترتیب میانگین اولویت و اندازه وظایف هستند. متغیر $l_t$ درجه محلی‌بودن وظایف، $\bar{v}_t$ میانگین سرعت خودروها، $\bar{\rho}_t$ میانگین بار واحدهای کنارجاده‌ای، و $\bar{\delta}_t$ میانگین زوال محبوبیت وظایف را نشان می‌دهند. این 9 مؤلفه وضعیت کامل سیستم را در هر لحظه توصیف می‌کنند.

تابع \lr{Q} با شبکه عصبی عمیق:
\begin{equation}
Q(s_t, a_t; \theta) = \text{DNN}(s_t, a_t; \theta)
\label{eq:q_function}
\end{equation}

در این معادله، $Q(s_t, a_t; \theta)$ تابع ارزش عمل-حالت است که ارزش انجام عمل $a_t$ در حالت $s_t$ را با پارامترهای شبکه $\theta$ محاسبه می‌کند. تابع \lr{DNN} شبکه عصبی عمیق با معماری $9 \rightarrow 64 \rightarrow 64 \rightarrow 5$ است که حالت $s_t$ را به عنوان ورودی دریافت کرده و \lr{Q-values} برای تمام 5 عمل ممکن ($a_t \in \{\text{\lr{noop}, \lr{evict\_lru}, \lr{evict\_random}, \lr{evict\_lfu}, \lr{shrink\_ttl}}\}$) را خروجی می‌دهد. پارامتر $\theta$ شامل تمام وزن‌ها و \lr{bias} های شبکه عصبی است.

به‌روزرسانی وزن‌ها با الگوریتم \lr{DQN}:
\begin{equation}
\theta_{t+1} = \theta_t + \alpha \left[ r_t + \gamma \max_{a'} Q(s_{t+1}, a'; \theta_t^-) - Q(s_t, a_t; \theta_t) \right] \nabla_{\theta_t} Q(s_t, a_t; \theta_t)
\label{eq:dqn_update}
\end{equation}

در این معادله، $\theta_{t+1}$ پارامترهای به‌روزرسانی شده شبکه در زمان $t+1$ هستند. پارامتر $\alpha$ نرخ یادگیری با مقدار 0.0001 است که سرعت به‌روزرسانی وزن‌ها را کنترل می‌کند. متغیر $r_t$ پاداش دریافتی در زمان $t$ و $\gamma = 0.95$ ضریب تخفیف است که اهمیت پاداش‌های آینده را تعیین می‌کند. عبارت $\max_{a'} Q(s_{t+1}, a'; \theta_t^-)$ حداکثر \lr{Q-value} در حالت بعدی با استفاده از شبکه هدف $\theta_t^-$ است که هر 100 گام به‌روزرسانی می‌شود. تفاوت $[r_t + \gamma \max_{a'} Q(s_{t+1}, a'; \theta_t^-) - Q(s_t, a_t; \theta_t)]$ خطای \lr{TD} (\lr{Temporal Difference}) را نمایش می‌دهد.

تابع پاداش چندهدفه:
\begin{equation}
R(s_t, a_t) = \alpha \cdot \Delta h_t + \beta \cdot \Delta \sigma_t - \gamma \cdot \Delta \bar{T}_t
\label{eq:reward_function}
\end{equation}

در این معادله، $R(s_t, a_t)$ تابع پاداش چندهدفه است که سه جنبه مهم عملکرد سیستم را در نظر می‌گیرد. متغیرهای $\Delta h_t = h_{t+1} - h_t$، $\Delta \sigma_t = \sigma_{t+1} - \sigma_t$، و $\Delta \bar{T}_t = \bar{T}_{t+1} - \bar{T}_t$ به ترتیب تغییرات در نرخ برخورد، نرخ موفقیت، و میانگین زمان پاسخ را نمایش می‌دهند. ضرایب وزن‌دهی $\alpha = 0.4$، $\beta = 0.4$، و $\gamma = 0.2$ اهمیت نسبی هر هدف را تعیین می‌کنند، به طوری که بهبود نرخ برخورد و نرخ موفقیت پاداش مثبت و افزایش زمان پاسخ پاداش منفی دریافت می‌کند. این ساختار تابع پاداش تضمین می‌کند که عامل یادگیری تقویتی به دنبال بهبود همزمان کارایی و کیفیت سرویس باشد.

\section{الگوریتم‌های کلیدی سیستم}

این بخش الگوریتم‌های اصلی سیستم حافظه نهان هوشمند را ارائه می‌دهد که شامل الگوریتم اصلی سیستم، الگوریتم یادگیری فدرال توزیع‌شده، و الگوریتم یادگیری تقویتی \lr{DQN} می‌باشد. این الگوریتم‌ها به صورت هماهنگ عمل کرده و عملکرد بهینه سیستم را تضمین می‌کنند.

الگوریتم اصلی سیستم هماهنگی کلی بین تمام اجزای سیستم را مدیریت می‌کند. این الگوریتم در هر گام زمانی وظایف جدید تولید کرده، آنها را پردازش می‌کند، و بر اساس فرکانس‌های تعریف شده مدل‌های یادگیری را به‌روزرسانی می‌کند. فرکانس آموزش محلی هر 10 گام و ثبت آمار هر 100 گام برای تعادل بین دقت و کارایی محاسباتی انتخاب شده‌اند.

الگوریتم یادگیری فدرال بر اساس \lr{FedAvg} پیاده‌سازی شده و شامل دو بخش اصلی است: آموزش محلی در هر واحد کنارجاده‌ای و تجمیع وزن‌دار در سرور مرکزی. تابع \lr{LocalUpdate} آموزش محلی را برای تعداد مشخصی دوره انجام داده و پارامترهای به‌روزرسانی شده را برمی‌گرداند. این الگوریتم امکان یادگیری توزیع‌شده را بدون نیاز به اشتراک‌گذاری داده‌های خام فراهم می‌کند.

الگوریتم \lr{DQN} شامل مکانیزم‌های کلیدی یادگیری تقویتی عمیق است: بازپخت تجربه، شبکه هدف، و استراتژی اکتشاف $\epsilon$-greedy. به‌روزرسانی دوره‌ای شبکه هدف پایداری آموزش را تضمین می‌کند و از واگرایی جلوگیری می‌کند. این الگوریتم قادر است سیاست‌های بهینه اخراج را در محیط پویای خودروی یاد بگیرد.

تعامل بین این سه الگوریتم منجر به ایجاد سیستمی می‌شود که قادر است به طور خودکار با تغییرات محیط تطبیق یافته و عملکرد بهینه را حفظ کند. یادگیری فدرال الگوهای محلی را شناسایی می‌کند، یادگیری تقویتی سیاست‌های سراسری را بهینه‌سازی می‌کند، و الگوریتم اصلی هماهنگی کلی را تضمین می‌کند.

\subsection{معماری کلی سیستم}

شکل \ref{fig:system_architecture} معماری کلی سیستم حافظه نهان هوشمند پیشنهادی را نشان می‌دهد که شامل واحدهای کنارجاده‌ای، کنترلر \lr{SDN}، سرور یادگیری فدرال، و عامل یادگیری تقویتی است.

\begin{figure}[h]
\centering
\includegraphics[width=0.9\textwidth]{system_architecture.png}
\caption{معماری کلی سیستم حافظه نهان هوشمند}
\label{fig:system_architecture}
\end{figure}

این معماری نشان می‌دهد که چگونه اجزای مختلف سیستم با یکدیگر تعامل می‌کنند. واحدهای کنارجاده‌ای داده‌های محلی را جمع‌آوری کرده و با کنترلر \lr{SDN} ارتباط برقرار می‌کنند. سرور یادگیری فدرال مدل‌های محلی را تجمیع کرده و عامل یادگیری تقویتی سیاست‌های بهینه را تعیین می‌کند.

\section{معیارهای ارزیابی عملکرد}

جهت ارزیابی جامع عملکرد سیستم، مجموعه‌ای از معیارهای کمی و کیفی تعریف شده که در جدول \ref{tab:performance_metrics} تشریح شده است. این معیارها پنج دسته اصلی را پوشش می‌دهند: عملکرد حافظه نهان، کیفیت سرویس، بهره‌وری منابع، شبکه و ارتباطات، و یادگیری و تطبیق.

\begin{table}[h]
\centering
\caption{معیارهای ارزیابی عملکرد سیستم}
\label{tab:performance_metrics}
\begin{tabular}{|p{3cm}|p{4cm}|p{6cm}|}
\hline
\rowcolor{lightgray}
\textbf{دسته‌بندی} & \textbf{معیار} & \textbf{فرمول محاسباتی} \\
\hline
\multirow{4}{3cm}{عملکرد حافظه نهان} & نرخ برخورد (\lr{Hit Rate}) & $HR = \frac{\text{تعداد برخوردها}}{\text{کل درخواست‌ها}} \times 100\%$ \\
\cline{2-3}
 & نرخ عدم برخورد (\lr{Miss Rate}) & $MR = 1 - HR$ \\
\cline{2-3}
 & میانگین زمان دسترسی & $\bar{T}_{access} = \frac{\sum_{i=1}^{n} T_i}{n}$ \\
\cline{2-3}
 & کارایی استفاده (\lr{Cache Utilization}) & $CU = \frac{\text{فضای اشغال‌شده}}{\text{کل ظرفیت}} \times 100\%$ \\
\hline
\multirow{3}{3cm}{کیفیت سرویس} & میانگین زمان پاسخ & $\bar{T}_{response} = \bar{T}_{process} + \bar{T}_{comm} + \bar{T}_{queue}$ \\
\cline{2-3}
 & نرخ موفقیت (\lr{Success Rate}) & $SR = \frac{\text{وظایف تکمیل‌شده}}{\text{کل وظایف}} \times 100\%$ \\
\cline{2-3}
 & نرخ نقض مهلت (\lr{Deadline Violation Rate}) & $DVR = \frac{\text{وظایف منقضی}}{\text{کل وظایف}} \times 100\%$ \\
\hline
\multirow{3}{3cm}{بهره‌وری منابع} & نرخ استفاده \lr{CPU} & $CPU_{util} = \frac{\text{زمان مشغول CPU}}{\text{کل زمان}} \times 100\%$ \\
\cline{2-3}
 & نرخ استفاده حافظه & $Mem_{util} = \frac{\text{حافظه استفاده‌شده}}{\text{کل حافظه}} \times 100\%$ \\
\cline{2-3}
 & توان مصرفی & $P_{total} = \sum_{r=1}^{R} P_{واحد\_کنارجاده‌ای_r} + P_{controller}$ \\
\hline
\multirow{2}{3cm}{شبکه و ارتباطات} & تأخیر انتشار & $D_{prop} = \frac{d}{c}$ (d: فاصله، c: سرعت نور) \\
\cline{2-3}
 & پهنای باند مؤثر & $BW_{eff} = \frac{\text{داده منتقل‌شده}}{\text{زمان انتقال}}$ \\
\hline
\multirow{2}{3cm}{یادگیری و تطبیق} & دقت پیش‌بینی \lr{FL} & $Acc_{FL} = \frac{TP + TN}{TP + TN + FP + FN}$ \\
\cline{2-3}
 & پاداش تجمعی \lr{RL} & $R_{cumulative} = \sum_{t=1}^{T} \gamma^{t-1} r_t$ \\
\hline
\end{tabular}
\end{table}

\subsection{تحلیل معیارهای عملکرد}

معیارهای عملکرد حافظه نهان شامل نرخ برخورد، نرخ عدم برخورد، میانگین زمان دسترسی، و کارایی استفاده می‌شوند. نرخ برخورد اصلی‌ترین معیار است که نشان‌دهنده درصد درخواست‌هایی است که از حافظه نهان محلی پاسخ داده می‌شوند. هدف افزایش این نرخ به بالای 80% است.

معیارهای کیفیت سرویس شامل میانگین زمان پاسخ، نرخ موفقیت، و نرخ نقض مهلت هستند. زمان پاسخ شامل سه جزء اصلی است: زمان پردازش، زمان ارتباط، و زمان انتظار در صف. هدف حفظ زمان پاسخ زیر 100 میلی‌ثانیه برای وظایف بحرانی است.

معیارهای بهره‌وری منابع نرخ استفاده \lr{CPU}، حافظه، و توان مصرفی را اندازه‌گیری می‌کنند. توان مصرفی کل شامل مصرف تمام واحدهای کنارجاده‌ای و کنترلر مرکزی است. هدف حفظ نرخ استفاده منابع بین 60-80% برای تعادل بین کارایی و پایداری است.

معیارهای شبکه و ارتباطات تأخیر انتشار و پهنای باند مؤثر را بررسی می‌کنند. تأخیر انتشار بر اساس فاصله فیزیکی و سرعت نور محاسبه می‌شود، در حالی که پهنای باند مؤثر نرخ واقعی انتقال داده را نشان می‌دهد.

معیارهای یادگیری و تطبیق دقت مدل یادگیری فدرال و پاداش تجمعی عامل یادگیری تقویتی را ارزیابی می‌کنند. دقت مدل \lr{FL} بر اساس ماتریس درهم‌ریختگی محاسبه شده و هدف دستیابی به دقت بالای 90% است. پاداش تجمعی \lr{RL} نشان‌دهنده عملکرد کلی عامل در طول زمان است.

\section{تنظیمات شبیه‌سازی}

ویژگی‌های شبیه‌سازی طوری تنظیم شده‌اند که نزدیک‌ترین حالت به سناریوهای واقعی را شبیه‌سازی کنند. جدول \ref{tab:simulation_params} این ویژگی‌ها را به تفصیل نشان می‌دهد. پارامترهای انتخاب شده بر اساس استانداردهای صنعتی و مطالعات تجربی در حوزه شبکه‌های خودرویی تعیین شده‌اند.

\begin{table}[h]
\centering
\caption{ویژگی‌های شبیه‌سازی}
\label{tab:simulation_params}
\begin{tabular}{|p{4cm}|p{3cm}|p{6cm}|}
\hline
\rowcolor{lightgray}
\textbf{ویژگی} & \textbf{مقدار} & \textbf{توضیحات} \\
\hline
زمان شبیه‌سازی & 3600 ثانیه & یک ساعت شبیه‌سازی پیوسته \\
\hline
منطقه شبیه‌سازی & $5 \times 5$ km² & شبکه شهری واقعی با 25 تقاطع \\
\hline
تعداد واحد کنارجاده‌ای & 5 واحد & پوشش کامل منطقه مورد نظر \\
\hline
تعداد خودرو & 50 خودرو & تراکم متوسط ترافیک شهری \\
\hline
برد ارتباطی & 300 متر & استاندارد \lr{DSRC} \\
\hline
\multirow{2}{4cm}{نرخ تولید وظیفه} & $\lambda = 0.3$ /ثانیه & فرآیند پواسون \\
 & (18 وظیفه/دقیقه) & توزیع تصادفی در زمان \\
\hline
\multirow{3}{4cm}{توزیع اولویت} & بالا: 30\% & وظایف ایمنی-حیاتی \\
 & متوسط: 50\% & وظایف کیفیت سرویس \\
 & پایین: 20\% & وظایف تفریحی \\
\hline
\multirow{4}{4cm}{اندازه وظایف} & کوچک: 10-100 \lr{KB} & درخواست‌های ساده \\
 & متوسط: 100\lr{KB}-1\lr{MB} & پردازش تصویر \\
 & بزرگ: 1-10 \lr{MB} & ویدیو و نقشه \\
 & خیلی بزرگ: >10 \lr{MB} & مدل‌های \lr{ML} \\
\hline
ظرفیت حافظه نهان & 100 ورودی/واحد & محدودیت فضای ذخیره‌سازی \\
\hline
\multirow{2}{4cm}{ویژگی‌های \lr{RL}} & $\epsilon$: 1.0 → 0.01 & کاهش تدریجی اکتشاف \\
 & $\gamma$: 0.95 & ضریب تخفیف آینده \\
\hline
\multirow{2}{4cm}{ویژگی‌های \lr{FL}} & \lr{Rounds}: 100 & دوره‌های آموزش فدرال \\
 & \lr{Local epochs}: 5 & تکرار محلی در هر واحد \\
\hline
\end{tabular}
\end{table}

\subsection{توجیه پارامترهای شبیه‌سازی}

زمان شبیه‌سازی 3600 ثانیه (یک ساعت) برای بررسی رفتار سیستم در بازه زمانی کافی انتخاب شده که امکان مشاهده تغییرات الگوهای ترافیک و عملکرد الگوریتم‌های یادگیری را فراهم می‌کند. منطقه شبیه‌سازی $5 \times 5$ کیلومتر مربع معادل یک ناحیه شهری متوسط است که شامل 25 تقاطع می‌شود.

تعداد 5 واحد کنارجاده‌ای با برد ارتباطی 300 متر پوشش کاملی از منطقه شبیه‌سازی فراهم می‌کند. این برد مطابق با استاندارد \lr{DSRC} (\lr{Dedicated Short Range Communications}) است که در شبکه‌های خودرویی استفاده می‌شود. تعداد 50 خودرو تراکم متوسط ترافیک شهری را شبیه‌سازی می‌کند.

نرخ تولید وظیفه با پارامتر $\lambda = 0.3$ بر ثانیه (معادل 18 وظیفه در دقیقه) بر اساس فرآیند پواسون تعریف شده که توزیع واقعی درخواست‌ها در شبکه‌های خودرویی را نمایش می‌دهد. توزیع اولویت وظایف شامل 30% وظایف بالا (ایمنی-حیاتی)، 50% متوسط (کیفیت سرویس)، و 20% پایین (تفریحی) است.

اندازه وظایف در چهار دسته تقسیم شده: کوچک (10-100 \lr{KB}) برای درخواست‌های ساده، متوسط (100\lr{KB}-1\lr{MB}) برای پردازش تصویر، بزرگ (1-10 \lr{MB}) برای ویدیو و نقشه، و خیلی بزرگ (>10 \lr{MB}) برای مدل‌های یادگیری ماشین.

پارامترهای یادگیری تقویتی شامل $\epsilon$ که از 1.0 به 0.01 کاهش می‌یابد تا تعادل بین اکتشاف و بهره‌برداری را تضمین کند، و $\gamma = 0.95$ به عنوان ضریب تخفیف که اهمیت پاداش‌های آینده را تعیین می‌کند. پارامترهای یادگیری فدرال شامل 100 دوره آموزش سراسری و 5 دوره آموزش محلی در هر واحد کنارجاده‌ای است.

\section{نتایج اولیه و تحلیل}

بر اساس شبیه‌سازی‌های اولیه انجام شده، نتایج امیدوارکننده‌ای به دست آمده که در جدول \ref{tab:preliminary_results} خلاصه شده است. این نتایج نشان‌دهنده عملکرد مطلوب سیستم پیشنهادی در مقایسه با روش‌های سنتی مدیریت حافظه نهان است.

\begin{table}[t!]
\centering
\caption{نتایج اولیه شبیه‌سازی}
\label{tab:preliminary_results}
\begin{tabular}{|p{4cm}|p{3cm}|p{6cm}|}
\hline
\rowcolor{lightgray}
\textbf{معیار} & \textbf{مقدار} & \textbf{مقایسه با روش‌های سنتی} \\
\hline
نرخ برخورد حافظه نهان & 87.3\% & بهبود 23\% نسبت به \lr{LRU} \\
\hline
میانگین زمان پاسخ & 45.2 ms & کاهش 31\% نسبت به \lr{Random} \\
\hline
نرخ موفقیت وظایف & 94.7\% & بهبود 18\% نسبت به \lr{FIFO} \\
\hline
بهره‌وری منابع & 91.4\% & بهبود 26\% در استفاده از \lr{CPU} \\
\hline
دقت پیش‌بینی \lr{FL} & 92.1\% & همگرایی در 15 دور آموزش \\
\hline
پاداش تجمعی \lr{RL} & +2847 & رشد مستمر و پایدار \\
\hline
مصرف انرژی & 78.6 W & کاهش 19\% نسبت به \lr{baseline} \\
\hline
تأخیر شبکه & 12.3 ms & در محدوده مطلوب \lr{real-time} \\
\hline
\end{tabular}
\end{table}

\subsection{تحلیل عملکرد یادگیری فدرال}

شکل \ref{fig:fl_convergence} نمودار همگرایی مدل‌های یادگیری فدرال را نشان می‌دهد. نتایج حاکی از آن است که سیستم پس از حدود 15 دور آموزش به حالت پایدار می‌رسد. دقت پیش‌بینی از 62\% در دور اول به 92.1\% در دور پایانی ارتقاء یافته است که نشان‌دهنده کارایی بالای الگوریتم \lr{FedAvg} در محیط توزیع‌شده است.

\begin{figure}[h]
\centering
\includegraphics[width=0.8\textwidth]{fl_convergence.png}
\caption{نمودار همگرایی یادگیری فدرال}
\label{fig:fl_convergence}
\end{figure}

مهم‌ترین نکته در این نمودار سرعت همگرایی بالا و پایداری مدل پس از دور 15 است. این امر نشان می‌دهد که الگوریتم قادر است در زمان کوتاهی به عملکرد مطلوب برسد و از نوسانات غیرضروری جلوگیری کند. تفاوت اندک بین دقت مدل سراسری و میانگین دقت محلی نیز حاکی از توزیع مناسب داده‌ها و عدم وجود \lr{data heterogeneity} شدید است.

\subsection{تحلیل عملکرد یادگیری تقویتی}

شکل \ref{fig:rl_performance} عملکرد عامل یادگیری تقویتی را در طول فرآیند آموزش نشان می‌دهد. عامل \lr{RL} پس از حدود 200 اپیزود آموزش، سیاست بهینه‌ای کسب کرده که منجر به بهبود قابل توجه در نرخ برخورد حافظه نهان و کاهش زمان پاسخ شده است.

\begin{figure}[h]
\centering
\includegraphics[width=0.8\textwidth]{rl_performance.png}
\caption{نمودار عملکرد یادگیری تقویتی}
\label{fig:rl_performance}
\end{figure}

پاداش تجمعی روند صعودی مستمری داشته و طی تمام دوره آموزش رشد قابل توجهی نشان می‌دهد. عامل \lr{RL} از پاداش منفی -420 در اپیزود اول به پاداش مثبت 1820 در اپیزود 300 رسیده است. تغییر از پاداش منفی به مثبت در حدود اپیزود 70 اتفاق افتاده که نشان‌دهنده شروع یادگیری مؤثر سیاست است. شیب ثابت رشد نشان می‌دهد که عامل به طور مداوم در حال بهبود عملکرد است و هنوز به نقطه اشباع نرسیده است.

\subsection{مقایسه با روش‌های موجود}

شکل \ref{fig:comparison} مقایسه جامعی از عملکرد سیستم پیشنهادی با روش‌های سنتی مدیریت حافظه نهان ارائه می‌دهد.

\begin{figure}[h]
\centering
\includegraphics[width=0.8\textwidth]{comparison.png}
\caption{مقایسه نرخ برخورد حافظه نهان}
\label{fig:comparison}
\end{figure}

نتایج نشان می‌دهد که سیستم پیشنهادی با نرخ برخورد 87.3\% عملکرد بهتری نسبت به تمام روش‌های سنتی دارد. بهبود 23\% نسبت به \lr{LRU} که بهترین روش سنتی محسوب می‌شود، نشان‌دهنده کارایی بالای ترکیب یادگیری فدرال و تقویتی است.

نتایج نشان می‌دهد که سیستم پیشنهادی در تمام معیارهای کلیدی بهبود قابل توجهی نسبت به روش‌های سنتی داشته است. بیشترین بهبود در زمان پاسخ (31%) و کمترین بهبود در همگرایی یادگیری فدرال (15%) مشاهده می‌شود که نشان‌دهنده تعادل مناسب بین سرعت و دقت سیستم است.

\subsection{تحلیل زمان پاسخ}

شکل \ref{fig:response_time} زمان پاسخ متوسط سیستم را در برابر تعداد وظایف مختلف نشان می‌دهد.

\begin{figure}[h]
\centering
\includegraphics[width=\textwidth]{response_time.png}
\caption{مقایسه زمان پاسخ متوسط}
\label{fig:response_time}
\end{figure}

نتایج نشان می‌دهد که سیستم پیشنهادی در تمام سطوح بار کاری عملکرد بهتری نسبت به روش‌های سنتی دارد. با افزایش تعداد وظایف، تفاوت عملکرد بیشتر می‌شود که نشان‌دهنده مقیاس‌پذیری مناسب سیستم است.

\subsubsection{مقایسه با الگوریتم RELIC}

برای ارزیابی دقیق‌تر عملکرد، سیستم پیشنهادی با الگوریتم \lr{RELIC} موجود در سیستم مقایسه شده است. \lr{RELIC} (Reliable and Low-Latency ML-based Trajectory Predictive Offloading) بر اساس بهینه‌سازی خطی دو مرحله‌ای عمل می‌کند:

\begin{itemize}
\item \textbf{مرحله اول:} تخصیص بهینه منابع محاسباتی با حداکثرسازی منابع اختصاص یافته
\item \textbf{مرحله دوم:} انتخاب سرور با حداقل‌سازی زمان تکمیل وظایف
\end{itemize}

\begin{figure}[h]
\centering
\includegraphics[width=\textwidth]{relic_comparison.png}
\caption{مقایسه عملکرد با الگوریتم RELIC}
\label{fig:relic_comparison}
\end{figure}

نتایج مقایسه با مهلت زمانی واقع‌گرایانه (\lr{0.5} ثانیه) نشان می‌دهد:

\begin{itemize}
\item \textbf{نرخ موفقیت:} \lr{RELIC} از \lr{85.2\%} به \lr{42.3\%} کاهش می‌یابد، \lr{FL-RL} از \lr{94.8\%} به \lr{78.9\%} (مقاومت بهتر)
\item \textbf{زمان تکمیل:} بهبود \lr{10.8\%} تا \lr{35.7\%} - نشان‌دهنده مزیت کش هوشمند در شرایط سخت
\end{itemize}

\textbf{تحلیل مقایسه با مهلت واقع‌گرایانه:} با استفاده از مهلت زمانی \lr{0.5} ثانیه (مناسب برای کاربردهای بلادرنگ خودرویی)، تفاوت عملکرد دو الگوریتم آشکار می‌شود. \lr{RELIC} به دلیل سربار بهینه‌سازی خطی، در شرایط سخت عملکرد ضعیفی دارد. سیستم \lr{FL-RL} با کش هوشمند و تصمیم‌گیری سریع، مقاومت بهتری نشان می‌دهد:

\begin{table}[h]
\centering
\caption{نتایج مقایسه با مهلت واقع‌گرایانه \lr{0.5} ثانیه}
\label{tab:realistic_comparison}
\begin{tabular}{|c|c|c|c|c|c|}
\hline
\textbf{تعداد وظایف} & \textbf{\lr{RELIC} نرخ موفقیت} & \textbf{\lr{FL-RL} نرخ موفقیت} & \textbf{\lr{RELIC} زمان (ثانیه)} & \textbf{\lr{FL-RL} زمان (ثانیه)} & \textbf{بهبود نرخ موفقیت} \\
\hline
250 & \lr{85.2\%} & \lr{94.8\%} & 0.485 & 0.312 & \lr{+9.6\%} \\
\hline
500 & \lr{72.4\%} & \lr{89.6\%} & 0.492 & 0.358 & \lr{+17.2\%} \\
\hline
750 & \lr{58.1\%} & \lr{84.2\%} & 0.498 & 0.401 & \lr{+26.1\%} \\
\hline
1000 & \lr{42.3\%} & \lr{78.9\%} & 0.499 & 0.445 & \lr{+36.6\%} \\
\hline
\end{tabular}
\end{table}

\textbf{دلایل برتری \lr{FL-RL} در زمان تکمیل:}

\begin{enumerate}
\item \textbf{کش هوشمند:} پیش‌قرارگیری وظایف محبوب کاهش زمان دسترسی
\item \textbf{عدم وابستگی به بهینه‌سازی:} \lr{RELIC} نیاز به حل مسائل خطی پیچیده دارد
\item \textbf{تصمیم‌گیری سریع:} شبکه عصبی آموزش‌دیده پاسخ فوری ارائه می‌دهد
\item \textbf{مقیاس‌پذیری:} عملکرد \lr{FL-RL} با افزایش بار کمتر تنزل می‌یابد
\end{enumerate}

دلیل برتری سیستم پیشنهادی عبارتند از:

\begin{enumerate}
\item \textbf{یادگیری تطبیقی:} قابلیت یادگیری از الگوهای گذشته و تطبیق با شرایط جدید
\item \textbf{پیش‌بینی هوشمند:} استفاده از یادگیری فدرال برای پیش‌بینی بهتر تقاضا
\item \textbf{تصمیم‌گیری بلادرنگ:} عدم وابستگی به حل مسائل بهینه‌سازی پیچیده
\item \textbf{مدیریت عدم قطعیت:} قابلیت مقابله با تغییرات غیرقابل پیش‌بینی محیط
\end{enumerate}

\subsection{تحلیل نرخ موفقیت}

شکل \ref{fig:success_rate} نرخ موفقیت سیستم در تکمیل وظایف را نشان می‌دهد.

\begin{figure}[h]
\centering
\includegraphics[width=\textwidth]{success_rate.png}
\caption{مقایسه نرخ موفقیت وظایف}
\label{fig:success_rate}
\end{figure}

سیستم پیشنهادی نرخ موفقیت بالای 94% را حفظ می‌کند که نشان‌دهنده قابلیت اطمینان بالای سیستم است. این در حالی است که روش‌های سنتی با افزایش بار کاری کاهش قابل توجهی در نرخ موفقیت نشان می‌دهند.

\subsection{تحلیل مصرف انرژی}

شکل \ref{fig:energy_consumption} مصرف انرژی سیستم را در سطوح مختلف بار کاری نشان می‌دهد.

\begin{figure}[h]
\centering
\includegraphics[width=\textwidth]{energy_consumption.png}
\caption{مقایسه مصرف انرژی}
\label{fig:energy_consumption}
\end{figure}

سیستم پیشنهادی مصرف انرژی کمتری نسبت به روش‌های سنتی دارد که نشان‌دهنده بهینه‌سازی مناسب منابع است. این امر برای شبکه‌های خودرویی که منابع انرژی محدود دارند بسیار مهم است.

\section{چالش‌ها و راه‌حل‌های پیشنهادی}

طی فرآیند پیاده‌سازی و آزمایش، چندین چالش مهم شناسایی و راه‌حل‌هایی برای آن‌ها ارائه شده که در جدول \ref{tab:challenges_solutions} تشریح شده است.

\begin{table}[h]
\centering
\caption{چالش‌ها و راه‌حل‌های پیشنهادی}
\label{tab:challenges_solutions}
\begin{tabular}{|p{4cm}|p{5cm}|p{4cm}|}
\hline
\rowcolor{lightgray}
\textbf{چالش} & \textbf{توضیح} & \textbf{راه‌حل پیشنهادی} \\
\hline
ناهمگونی داده‌ها & توزیع غیریکنواخت داده‌ها در واحدهای مختلف & استفاده از تکنیک‌های \lr{data augmentation} و وزن‌دهی تطبیقی \\
\hline
تأخیر ارتباطات & تأخیر در انتقال مدل‌های \lr{FL} & فشرده‌سازی مدل و انتقال تدریجی پارامترها \\
\hline
پیچیدگی محاسباتی & بار محاسباتی بالای \lr{DQN} & استفاده از شبکه‌های کوچک‌تر و تکنیک‌های \lr{pruning} \\
\hline
مقیاس‌پذیری & افزایش تعداد واحدها & معماری سلسله‌مراتبی و خوشه‌بندی واحدها \\
\hline
امنیت و حریم خصوصی & حفاظت از داده‌های حساس & رمزنگاری همومورفیک و \lr{differential privacy} \\
\hline
تطبیق با تغییرات & تغییرات سریع الگوهای ترافیک & مکانیزم تطبیق آنلاین و به‌روزرسانی مداوم \\
\hline
\end{tabular}
\end{table}

\subsection{تحلیل چالش‌ها}

ناهمگونی داده‌ها یکی از اصلی‌ترین چالش‌های یادگیری فدرال است که می‌تواند منجر به همگرایی کند یا کیفیت مدل سراسری را کاهش دهد. راه‌حل پیشنهادی شامل استفاده از تکنیک‌های تقویت داده و وزن‌دهی تطبیقی بر اساس کیفیت داده‌های محلی است.

تأخیر ارتباطات در شبکه‌های خودرویی به دلیل تحرک بالا و تغییرات توپولوژی اجتناب‌ناپذیر است. فشرده‌سازی مدل‌ها و انتقال تدریجی پارامترها می‌تواند این تأخیر را کاهش دهد.

پیچیدگی محاسباتی \lr{DQN} در محیط‌های منابع محدود چالش‌برانگیز است. استفاده از معماری‌های کوچک‌تر و تکنیک‌های هرس کردن شبکه می‌تواند این مشکل را حل کند.

مقیاس‌پذیری سیستم برای شبکه‌های بزرگ نیازمند معماری سلسله‌مراتبی است که امکان مدیریت کارآمد تعداد زیادی واحد کنارجاده‌ای را فراهم کند.

\section{فعالیت‌های پیش‌رو}

این بخش برنامه‌ریزی جامعی از فعالیت‌های آینده پژوهش ارائه می‌دهد که شامل مقایسه‌های تفصیلی با کارهای قبلی، تدوین مقالات علمی، و توسعه ایده‌های نوآورانه است.

\subsection{مقایسه جامع با روش‌های موجود}

برای تکمیل ارزیابی سیستم پیشنهادی، مقایسه‌های تفصیلی با روش‌های پیشرفته موجود در ادبیات ضروری است:

\subsubsection{مقایسه با روش‌های یادگیری فدرال}

\begin{table}[h]
\centering
\caption{مقایسه با روش‌های یادگیری فدرال موجود}
\label{tab:fl_comparison}
\begin{tabular}{|p{3cm}|p{3cm}|p{3cm}|p{4cm}|}
\hline
\textbf{روش} & \textbf{دقت (\%)} & \textbf{زمان همگرایی} & \textbf{ویژگی‌های کلیدی} \\
\hline
\lr{FedAvg} کلاسیک & 89.2 & 25 دور & تجمیع ساده وزن‌ها \\
\hline
\lr{FedProx} & 90.1 & 22 دور & تنظیم‌سازی محلی \\
\hline
\lr{SCAFFOLD} & 90.8 & 20 دور & کنترل واریانس \\
\hline
روش پیشنهادی & 92.1 & 15 دور & تلفیق با \lr{RL} \\
\hline
\end{tabular}
\end{table}

\subsubsection{مقایسه با روش‌های یادگیری تقویتی}

\begin{table}[h]
\centering
\caption{مقایسه با الگوریتم‌های یادگیری تقویتی}
\label{tab:rl_comparison}
\begin{tabular}{|p{3cm}|p{3cm}|p{3cm}|p{4cm}|}
\hline
\textbf{الگوریتم} & \textbf{پاداش نهایی} & \textbf{زمان همگرایی} & \textbf{پیچیدگی محاسباتی} \\
\hline
\lr{Q-Learning} & 1420 & 450 اپیزود & $O(|S| \times |A|)$ \\
\hline
\lr{SARSA} & 1380 & 420 اپیزود & $O(|S| \times |A|)$ \\
\hline
\lr{Double DQN} & 1650 & 350 اپیزود & $O(n \times m)$ \\
\hline
روش پیشنهادی & 1820 & 300 اپیزود & $O(n \times m)$ \\
\hline
\end{tabular}
\end{table}

\subsubsection{مقایسه با سیستم‌های حافظه نهان هوشمند}

مطالعات اخیر در زمینه حافظه نهان هوشمند شامل:

\begin{itemize}
\item \textbf{\lr{DeepCache}}: استفاده از شبکه‌های عمیق برای پیش‌بینی الگوهای دسترسی
\item \textbf{\lr{MLCache}}: بهره‌گیری از یادگیری ماشین برای بهینه‌سازی جایگذاری
\item \textbf{\lr{EdgeAI-Cache}}: تلفیق محاسبات لبه با هوش مصنوعی
\item \textbf{\lr{FedEdge}}: یادگیری فدرال برای محاسبات لبه
\end{itemize}

\subsection{برنامه تدوین مقالات علمی}

بر اساس نتایج به‌دست آمده، برنامه‌ای برای انتشار مقالات علمی در مجلات معتبر تدوین شده است:

\subsubsection{مقاله اول: مبانی نظری}

\textbf{عنوان پیشنهادی:} \lr{"Federated Learning-Enhanced Intelligent Caching for Vehicular Networks: A Deep Reinforcement Learning Approach"}

\textbf{مجله هدف:} \lr{IEEE Transactions on Vehicular Technology}

\textbf{محتوای کلیدی:}
\begin{itemize}
\item معرفی معماری سیستم پیشنهادی
\item تحلیل نظری الگوریتم‌های \lr{FL} و \lr{RL}
\item اثبات همگرایی و پایداری سیستم
\item مقایسه پیچیدگی محاسباتی
\end{itemize}

\textbf{زمان‌بندی:} 3 ماه آینده

\subsubsection{مقاله دوم: ارزیابی عملکرد}

\textbf{عنوان پیشنهادی:} \lr{"Performance Evaluation of Hybrid FL-RL Caching System in Dynamic Vehicular Environments"}

\textbf{مجله هدف:} \lr{Computer Networks} یا \lr{IEEE Network}

\textbf{محتوای کلیدی:}
\begin{itemize}
\item نتایج شبیه‌سازی جامع
\item مقایسه با روش‌های پیشرفته
\item تحلیل حساسیت پارامترها
\item بررسی مقیاس‌پذیری
\end{itemize}

\textbf{زمان‌بندی:} 4-6 ماه آینده

\subsubsection{مقاله سوم: پیاده‌سازی و آزمایش}

\textbf{عنوان پیشنهادی:} \lr{"Real-world Implementation and Testbed Evaluation of Intelligent Vehicular Caching System"}

\textbf{مجله هدف:} \lr{IEEE Transactions on Mobile Computing}

\textbf{محتوای کلیدی:}
\begin{itemize}
\item پیاده‌سازی روی پلتفرم واقعی
\item آزمایش در محیط‌های مختلف
\item تحلیل مصرف منابع
\item راهنمای استقرار
\end{itemize}

\textbf{زمان‌بندی:} 6-9 ماه آینده

\subsection{ایده‌های نوآورانه برای توسعه}

\subsubsection{تلفیق با تکنولوژی‌های نوظهور}

\paragraph{\lr{5G/6G Integration}}
ادغام سیستم با شبکه‌های \lr{5G} و \lr{6G} برای:
\begin{itemize}
\item استفاده از \lr{Network Slicing} برای تخصیص منابع
\item بهره‌گیری از \lr{Ultra-Low Latency} برای کاربردهای بحرانی
\item پشتیبانی از \lr{Massive IoT} در شبکه‌های خودرویی
\end{itemize}

\paragraph{تلفیق با \lr{Blockchain}}
استفاده از بلاک‌چین برای:
\begin{itemize}
\item تضمین امنیت و حریم خصوصی داده‌ها
\item ایجاد سیستم پاداش غیرمتمرکز
\item اعتبارسنجی تراکنش‌های حافظه نهان
\end{itemize}

\paragraph{هوش مصنوعی تولیدی (\lr{Generative AI})}
بهره‌گیری از \lr{GPT} و مدل‌های زبان بزرگ برای:
\begin{itemize}
\item تولید سناریوهای آزمایش هوشمند
\item بهینه‌سازی خودکار پارامترها
\item تحلیل پیشرفته الگوهای ترافیک
\end{itemize}

\subsubsection{گسترش دامنه کاربرد}

\paragraph{شهرهای هوشمند}
توسعه سیستم برای:
\begin{itemize}
\item مدیریت ترافیک شهری
\item بهینه‌سازی مصرف انرژی
\item خدمات اضطراری هوشمند
\end{itemize}

\paragraph{اینترنت اشیاء صنعتی (\lr{IIoT})}
تطبیق سیستم برای:
\begin{itemize}
\item کارخانه‌های هوشمند
\item زنجیره تأمین خودکار
\item نگهداری پیش‌بینانه
\end{itemize}

\paragraph{محاسبات لبه چندمقیاسه}
توسعه معماری سلسله‌مراتبی شامل:
\begin{itemize}
\item \lr{Nano Edge}: سطح دستگاه
\item \lr{Micro Edge}: سطح محلی
\item \lr{Macro Edge}: سطح منطقه‌ای
\item \lr{Cloud Edge}: سطح ملی/بین‌المللی
\end{itemize}

\subsubsection{چالش‌های تحقیقاتی آینده}

\paragraph{بهینه‌سازی چندهدفه}
توسعه الگوریتم‌هایی که همزمان بهینه‌سازی کنند:
\begin{itemize}
\item کیفیت سرویس
\item مصرف انرژی
\item هزینه عملیاتی
\item امنیت سیستم
\end{itemize}

\paragraph{یادگیری انتقالی (\lr{Transfer Learning})}
امکان انتقال دانش بین:
\begin{itemize}
\item محیط‌های مختلف (شهری، جاده‌ای، بزرگراهی)
\item شرایط آب‌وهوایی متفاوت
\item الگوهای ترافیکی متنوع
\end{itemize}

\paragraph{تحلیل عدم قطعیت}
مدل‌سازی و مدیریت:
\begin{itemize}
\item عدم قطعیت در پیش‌بینی‌ها
\item تغییرات غیرقابل پیش‌بینی محیط
\item خرابی‌های تصادفی سیستم
\end{itemize}

\subsection{زمان‌بندی کلی پروژه بر اساس جدول تفصیلی}

بر اساس جدول زمان‌بندی تفصیلی ارائه شده، فعالیت‌های باقی‌مانده و آینده به شرح زیر برنامه‌ریزی شده‌اند:

\subsubsection{فاز 4: تکمیل و ارسال مقاله FL-RL کش (مهر 1404)}

\begin{itemize}
\item \textbf{ارسال مقاله FL-RL کش:} تکمیل \lr{5\%} باقی‌مانده و ارسال به \lr{IEEE Internet of Things Journal}
\item \textbf{مطالعه برای ایده جدید:} شناسایی چالش‌های جدید و طراحی مفهومی راه‌حل نوآورانه (1 ماه)
\item \textbf{پیاده‌سازی ایده جدید:} توسعه الگوریتم‌های جدید و محیط شبیه‌سازی (4 ماه)
\end{itemize}

\subsubsection{فاز 5: پژوهش و تدوین مقاله دوم (اردیبهشت-تیر 1405)}

\begin{itemize}
\item \textbf{تدوین مقاله دوم:} نگارش مقاله کامل انگلیسی با تمرکز بر عملکرد سیستم (3 ماه)
\item \textbf{هدف انتشار:} \lr{IEEE Transactions on Vehicular Technology}
\item \textbf{تکمیل رساله دکتری:} تلفیق تمام فصول و آماده‌سازی برای دفاع (3 ماه)
\end{itemize}

\begin{table}[h]
\centering
\caption{برنامه زمان‌بندی فعالیت‌های پیش‌رو بر اساس جدول تفصیلی}
\label{tab:timeline_updated}
\begin{tabular}{|p{2.5cm}|p{4cm}|p{2.5cm}|p{3cm}|}
\hline
\textbf{فاز} & \textbf{فعالیت} & \textbf{زمان} & \textbf{وضعیت} \\
\hline
\textbf{فاز 4} & ارسال مقاله FL-RL کش & مهر 1404 & \textcolor{blue}{برنامه‌ریزی شده} \\
\hline
& مطالعه ایده جدید & شهریور 1404 & \textcolor{blue}{برنامه‌ریزی شده} \\
\hline
& پیاده‌سازی & دی 1404 & \textcolor{blue}{برنامه‌ریزی شده} \\
\hline
\textbf{فاز 5} & تدوین مقاله دوم & اردیبهشت 1405 & \textcolor{blue}{برنامه‌ریزی شده} \\
\hline
& تکمیل رساله دکتری & مرداد 1405 & \textcolor{blue}{برنامه‌ریزی شده} \\
\hline
\end{tabular}
\end{table}

\subsubsection{اهداف کلیدی آینده}

\begin{enumerate}
\item \textbf{انتشار علمی:} تکمیل حداقل 3 مقاله در مجلات معتبر \lr{ISI}
\item \textbf{نوآوری تکنولوژیکی:} توسعه سیستم کش هوشمند قابل تجاری‌سازی  
\item \textbf{تکمیل رساله:} آماده‌سازی برای دفاع دکتری تا مهر 1405
\item \textbf{تأثیر علمی:} ایجاد استاندارد جدید در \lr{Intelligent Caching}
\end{enumerate}

\subsection{منابع مورد نیاز}

برای تحقق اهداف فوق، منابع زیر ضروری است:

\subsubsection{منابع سخت‌افزاری}
\begin{itemize}
\item سرور محاسباتی با \lr{GPU} برای شبیه‌سازی‌های سنگین
\item تجهیزات شبکه برای آزمایش‌های عملی
\item دستگاه‌های \lr{IoT} و سنسورها
\end{itemize}

\subsubsection{منابع نرم‌افزاری}
\begin{itemize}
\item مجوز نرم‌افزارهای شبیه‌سازی (\lr{SUMO, NS-3, MATLAB})
\item پلتفرم‌های یادگیری ماشین (\lr{TensorFlow, PyTorch})
\item ابزارهای تحلیل داده (\lr{Python, R})
\end{itemize}

\subsubsection{منابع انسانی}
\begin{itemize}
\item همکاری با محققان متخصص در حوزه \lr{FL} و \lr{RL}
\item مشاوره با متخصصان صنعت خودرو
\item همکاری بین‌المللی با دانشگاه‌های معتبر
\end{itemize}

این برنامه جامع مسیر روشنی برای ادامه تحقیق و دستیابی به نتایج قابل انتشار در مجلات معتبر بین‌المللی فراهم می‌کند.

\section{جمع‌بندی}

این پژوهش سیستم حافظه نهان هوشمند برای شبکه‌های خودرویی را با تلفیق یادگیری فدرال و یادگیری تقویتی عمیق ارائه داده است. نتایج به‌دست آمده نشان می‌دهد که رویکرد پیشنهادی قادر است چالش‌های اصلی شبکه‌های خودرویی شامل تحرک بالا، تغییرات توپولوژی، و حفظ حریم خصوصی را به طور مؤثر حل کند.

\subsection{دستاوردهای کلیدی}

مهم‌ترین دستاوردهای این پژوهش عبارتند از:

\begin{itemize}
\item \textbf{انتشارات علمی:} ارسال 2 مقاله به مجلات معتبر بین‌المللی شامل:
  \begin{itemize}
  \item \lr{RELIC}: سیستم پیش‌بینی مسیر و تخصیص منابع خودروهای پارک شده
  \item \lr{FogZoneSim}: شبیه‌ساز منطقه‌ای برای مدیریت منابع در شبکه‌های IoT-Fog
  \end{itemize}
\item \textbf{معماری نوآورانه:} طراحی سیستم یکپارچه FL-RL که قابلیت‌های یادگیری توزیع‌شده و تصمیم‌گیری هوشمند را ترکیب می‌کند
\item \textbf{عملکرد برتر:} بهبود 23\% در نرخ برخورد حافظه نهان و کاهش 31\% در زمان پاسخ نسبت به روش‌های سنتی
\item \textbf{مقیاس‌پذیری:} قابلیت مدیریت شبکه‌های بزرگ با تعداد زیادی واحد کنارجاده‌ای و خودرو
\item \textbf{حفظ حریم خصوصی:} استفاده از یادگیری فدرال برای محافظت از داده‌های حساس خودروها
\item \textbf{کاهش اجرای مجدد وظایف:} بهبود نرخ موفقیت از طریق کش هوشمند وظایف
\end{itemize}

\subsection{تأثیرات علمی و عملی}

این پژوهش تأثیرات مهمی در حوزه‌های مختلف دارد:

\subsubsection{تأثیرات علمی}
\begin{itemize}
\item ارائه چارچوب جدید برای تلفیق FL و RL در شبکه‌های خودرویی
\item توسعه الگوریتم‌های بهینه‌سازی چندهدفه برای مدیریت حافظه نهان
\item مشارکت در پیشرفت دانش در حوزه محاسبات لبه هوشمند
\end{itemize}

\subsubsection{تأثیرات عملی}
\begin{itemize}
\item بهبود تجربه کاربری در خدمات خودرویی
\item کاهش مصرف پهنای باند و انرژی در شبکه‌های حمل‌ونقل
\item امکان استقرار خدمات هوشمند در شهرهای آینده
\end{itemize}

\subsection{جدول زمان‌بندی فعالیت‌ها}

برنامه زمان‌بندی تفصیلی برای تکمیل پروژه و انتشار نتایج در جدول \ref{tab:detailed_timeline} ارائه شده است:

\begin{table}[h]
\centering
\caption{جدول زمان‌بندی تفصیلی فعالیت‌های پروژه}
\label{tab:detailed_timeline}
\resizebox{\textwidth}{!}{
\begin{tabular}{|c|p{4cm}|p{3cm}|p{2cm}|p{3cm}|p{2cm}|}
\hline
\textbf{فاز} & \textbf{فعالیت} & \textbf{مدت زمان} & \textbf{شروع} & \textbf{خروجی} & \textbf{وضعیت} \\
\hline
\multirow{3}{*}{\textbf{فاز 1}} & تدوین مقاله RELIC & 3 ماه & مهر 1403 & مقاله RELIC & \textcolor{green}{تکمیل} \\
\cline{2-6}
& ارسال مقاله RELIC & 1 هفته & دی 1403 & ارسال & \textcolor{green}{تکمیل} \\
\cline{2-6}
& شروع پروژه FL-RL کش & 1 ماه & دی 1403 & طرح کلی & \textcolor{green}{تکمیل} \\
\hline
\multirow{4}{*}{\textbf{فاز 2}} & تدوین مقاله FogZoneSim & 4 ماه & اسفند 1403 & مقاله شبیه‌ساز & \textcolor{green}{تکمیل} \\
\cline{2-6}
& ارسال مقاله FogZoneSim & 1 هفته & اردیبهشت 1404 & ارسال & \textcolor{green}{تکمیل} \\
\cline{2-6}
& پیاده‌سازی الگوریتم FL & 2 ماه & بهمن 1403 & کد FL & \textcolor{green}{تکمیل} \\
\cline{2-6}
& پیاده‌سازی الگوریتم RL & 2 ماه & اسفند 1403 & کد DQN & \textcolor{green}{تکمیل} \\
\hline
\multirow{3}{*}{\textbf{فاز 3}} & تلفیق سیستم FL-RL & 1 ماه & خرداد 1404 & سیستم یکپارچه & \textcolor{green}{تکمیل} \\
\cline{2-6}
& شبیه‌سازی و تولید نمودارها & 2 ماه & تیر 1404 & نتایج و تصاویر & \textcolor{green}{تکمیل} \\
\cline{2-6}
& تدوین مقاله FL-RL کش & 2 ماه & مرداد 1404 & مقاله انگلیسی & \textcolor{orange}{در حال انجام} \\
\hline
\multirow{3}{*}{\textbf{فاز 4}} & ارسال مقاله FL-RL کش & 1 هفته & مهر 1404 & ارسال & \textcolor{blue}{برنامه‌ریزی شده} \\
\cline{2-6}
& تحلیل حساسیت پارامترها & 1 ماه & شهریور 1404 & نتایج تحلیل & \textcolor{blue}{برنامه‌ریزی شده} \\
\cline{2-6}
& بررسی مقیاس‌پذیری & 1 ماه & مهر 1404 & آزمایش‌های بزرگ & \textcolor{blue}{برنامه‌ریزی شده} \\
\hline
\multirow{3}{*}{\textbf{فاز 5}} & تدوین مقاله دوم & 1 ماه & آبان 1404 & مقاله عملکرد & \textcolor{blue}{برنامه‌ریزی شده} \\
\cline{2-6}
& پیاده‌سازی نمونه اولیه & 2 ماه & دی 1404 & سیستم واقعی & \textcolor{blue}{برنامه‌ریزی شده} \\
\cline{2-6}
& آزمایش‌های عملی & 1 ماه & اسفند 1404 & نتایج واقعی & \textcolor{blue}{برنامه‌ریزی شده} \\
\hline
\multirow{2}{*}{\textbf{فاز 6}} & تدوین مقاله سوم & 1 ماه & فروردین 1405 & مقاله پیاده‌سازی & \textcolor{blue}{برنامه‌ریزی شده} \\
\cline{2-6}
& تکمیل رساله دکتری & 2 ماه & اردیبهشت 1405 & رساله نهایی & \textcolor{blue}{برنامه‌ریزی شده} \\
\hline
\end{tabular}
}
\end{table}

\subsection{جدول مقالات پذیرفته‌شده، در حال داوری و در دست تهیه}

وضعیت مقالات مرتبط با این پژوهش در جدول \ref{tab:papers_status} خلاصه شده است:

\begin{table}[h]
\centering
\caption{وضعیت مقالات پژوهشی}
\label{tab:papers_status}
\resizebox{\textwidth}{!}{
\begin{tabular}{|c|p{5cm}|p{3cm}|p{2cm}|p{2cm}|p{2cm}|}
\hline
\textbf{ردیف} & \textbf{عنوان مقاله} & \textbf{مجله/کنفرانس هدف} & \textbf{وضعیت} & \textbf{تاریخ ارسال} & \textbf{درصد تکمیل} \\
\hline
1 & \lr{RELIC: Reliable and Low-Latency ML-based Trajectory Predictive Offloading for SD-VANETs with Efficient Parked Vehicle Resource Utilization} & \lr{IEEE Transactions on Vehicular Technology} & \textcolor{red}{ارسال شده} & دی 1403 & 100\% \\
\hline
2 & \lr{Federated Learning-Enhanced Intelligent Caching for Vehicular Networks: A Deep Reinforcement Learning Approach} & \lr{Computer Networks} & در دست تهیه & مهر 1404 & 95\% \\
\hline
3 & \lr{Performance Evaluation of Hybrid FL-RL Caching System in Dynamic Vehicular Environments} & \lr{IEEE Transactions on Mobile Computing} & برنامه‌ریزی شده & آبان 1404 & 60\% \\
\hline
4 & \lr{Real-world Implementation and Testbed Evaluation of Intelligent Vehicular Caching System} & \lr{IEEE Network} & برنامه‌ریزی شده & دی 1404 & 30\% \\
\hline
5 & \lr{Privacy-Preserving Federated Learning for Edge Caching in VANETs} & \lr{IEEE Internet of Things Journal} & ایده اولیه & اسفند 1404 & 15\% \\
\hline
\multicolumn{6}{|c|}{\textbf{مقالات کنفرانس}} \\
\hline
6 & \lr{FogZoneSim: Zone-Based Simulator for Large-Scale Resource Management in IoT-Fog Networks} & \lr{IEEE INFOCOM 2024} & \textcolor{red}{ارسال شده} & اردیبهشت 1404 & 100\% \\
\hline
7 & \lr{Preliminary Results of FL-RL Caching in Vehicular Networks} & \lr{IEEE VTC 2024} & برنامه‌ریزی شده & آبان 1404 & 70\% \\
\hline
\multicolumn{6}{|c|}{\textbf{مقالات داخلی}} \\
\hline
8 & سیستم حافظه نهان هوشمند برای شبکه‌های خودرویی با رویکرد یادگیری فدرال & مجله مهندسی برق دانشگاه تبریز & در دست تهیه & مرداد 1404 & 80\% \\
\hline
9 & کاربرد یادگیری تقویتی در بهینه‌سازی شبکه‌های خودرویی & نشریه مهندسی کامپیوتر ایران & برنامه‌ریزی شده & مهر 1404 & 40\% \\
\hline
\end{tabular}
}
\end{table}

\subsection{نقشه راه تحقیقاتی}

بر اساس پیشرفت‌های انجام شده، نقشه راه کلی پروژه به شرح زیر است:

\subsubsection{مرحله کوتاه‌مدت (6 ماه آینده)}
\begin{itemize}
\item تکمیل مقایسه‌های جامع با روش‌های موجود
\item ارسال اولین مقاله به مجله معتبر بین‌المللی
\item شرکت در کنفرانس‌های علمی برای ارائه نتایج اولیه
\item شروع همکاری با گروه‌های تحقیقاتی بین‌المللی
\end{itemize}

\subsubsection{مرحله میان‌مدت (6-12 ماه)}
\begin{itemize}
\item پیاده‌سازی نمونه اولیه سیستم
\item انجام آزمایش‌های عملی در محیط کنترل شده
\item ارسال مقالات دوم و سوم
\item تکمیل فصل‌های رساله دکتری
\end{itemize}

\subsubsection{مرحله بلندمدت (12-18 ماه)}
\begin{itemize}
\item استقرار سیستم در محیط واقعی
\item همکاری با صنعت برای تجاری‌سازی
\item توسعه نسخه‌های پیشرفته سیستم
\item دفاع از رساله دکتری
\end{itemize}

\subsection{شاخص‌های موفقیت}

معیارهای ارزیابی موفقیت پروژه عبارتند از:

\begin{table}[h]
\centering
\caption{شاخص‌های موفقیت پروژه}
\label{tab:success_metrics}
\begin{tabular}{|p{4cm}|p{3cm}|p{3cm}|p{3cm}|}
\hline
\textbf{شاخص} & \textbf{هدف} & \textbf{وضعیت فعلی} & \textbf{پیش‌بینی نهایی} \\
\hline
تعداد مقالات ISI & 3 مقاله & 2 ارسال شده & 4-5 مقاله \\
\hline
تعداد استنادات & 50+ استناد & در انتظار پذیرش & 30-50 استناد \\
\hline
ثبت اختراع & 1 اختراع & در حال بررسی & 1-2 اختراع \\
\hline
همکاری صنعتی & 1 قرارداد & مذاکره اولیه & 1-2 قرارداد \\
\hline
نمونه اولیه & سیستم کامل & 85\% تکمیل & سیستم عملیاتی \\
\hline
\end{tabular}
\end{table}

\subsection{نتیجه‌گیری نهایی}

این پژوهش با موفقیت نشان داده است که تلفیق یادگیری فدرال و یادگیری تقویتی عمیق می‌تواند راه‌حل مؤثری برای چالش‌های حافظه نهان در شبکه‌های خودرویی ارائه دهد. نتایج به‌دست آمده نه تنها از نظر علمی قابل توجه است، بلکه پتانسیل کاربردهای عملی گسترده‌ای در صنعت حمل‌ونقل هوشمند دارد.

با پیروی از برنامه زمان‌بندی ارائه شده و تکمیل فعالیت‌های پیش‌رو، این پروژه قادر خواهد بود تأثیر مهمی در پیشرفت دانش و فناوری شبکه‌های خودرویی داشته باشد و زمینه را برای تحقیقات آینده در این حوزه فراهم کند.

\end{document}
